% supplement.tex -- Supplementary Material. SOURCE OF TRUTH from 2026-10-15
% (docs/manuscript/SUPPLEMENT.md is a frozen snapshot). Build: ./build.sh (tectonic).
% Cross-references to the main text are by section name/number as printed in main.pdf ("M2.4", "main Table 2").
\documentclass[11pt,a4paper]{article}
% macros.tex -- shared preamble for main.tex and supplement.tex.
% Load with % macros.tex -- shared preamble for main.tex and supplement.tex.
% Load with % macros.tex -- shared preamble for main.tex and supplement.tex.
% Load with \input{macros} directly after \documentclass[11pt,a4paper]{article}.
% Built with tectonic (XeTeX engine); also compiles with pdflatex.

% ---------------------------------------------------------------- fonts
\usepackage{iftex}
\ifPDFTeX
  \usepackage[T1]{fontenc}
  \usepackage{lmodern}
\else
  % XeTeX/LuaTeX: fontspec's default text font is Latin Modern (OpenType lmodern)
  \usepackage{fontspec}
\fi
\usepackage{textcomp}
\usepackage{microtype}
\usepackage[raggedright]{titlesec}

% ---------------------------------------------------------------- page layout
\usepackage[a4paper,margin=2.5cm,headheight=14pt]{geometry}
\linespread{1.3}
\usepackage{fancyhdr}
\pagestyle{fancy}
\fancyhf{}
\renewcommand{\headrulewidth}{0pt}
\fancyfoot[C]{\thepage}
\setlength{\parindent}{0pt}
\setlength{\parskip}{0.6\baselineskip}
\setlength{\emergencystretch}{3em}
\usepackage[running]{lineno} % continuous line numbers for reviewer comments
\renewcommand{\linenumberfont}{\normalfont\scriptsize\sffamily\color{gray}}

% ---------------------------------------------------------------- tables, figures
\usepackage{graphicx}
\graphicspath{{../docs/manuscript/figures/}} % single source: figures are not copied
\usepackage{booktabs}
\usepackage{array}
\usepackage{tabularx}
\usepackage{threeparttable}
\usepackage{pdflscape}
\usepackage{float}
\usepackage{placeins}
\usepackage{needspace}
\usepackage[titles]{tocloft}
\setlength{\cftsecnumwidth}{2.6em}
\setlength{\cftsubsecindent}{2.6em}
\setlength{\cftsubsecnumwidth}{3.4em}
\setlength{\cftbeforesecskip}{2pt}
\usepackage[font=small,labelfont=bf,labelsep=period,justification=justified,singlelinecheck=false]{caption}
\usepackage{siunitx}
\newcolumntype{L}{>{\raggedright\arraybackslash}X}
\newcolumntype{R}[1]{>{\raggedright\arraybackslash}p{#1}}
% U: like R, but \_ may break (plain-text variable names)
\newcolumntype{U}[1]{>{\raggedright\arraybackslash\renewcommand{\_}{\textunderscore\allowbreak}}p{#1}}
\setlength{\tabcolsep}{4pt}
\renewcommand{\arraystretch}{1.1}
\usepackage{etoolbox}
% floats (tables, figure legends) are single-spaced; running text keeps \linespread{1.3}
\AtBeginEnvironment{table}{\linespread{1}\selectfont}
\AtBeginEnvironment{figure}{\linespread{1}\selectfont}

% ---------------------------------------------------------------- colour, placeholders
\usepackage{xcolor}
% Visible placeholders: [CHECK ...] citations and [PENDING ...] author decisions.
\newcommand{\flag}[1]{\textcolor{red}{#1}}
\newcommand{\todo}[1]{\textcolor{red}{[#1]}}

% ---------------------------------------------------------------- inline code / file paths
% \code{...}: typewriter text that may break after "/", "_" and "," (robust, safe in captions).
% Inside \code, escape _ % & # { } as usual (\_ etc.).
\ExplSyntaxOn
\tl_new:N \l__paper_code_tl
\NewDocumentCommand{\code}{m}
  {
    \group_begin:
    \ttfamily\small
    \tl_set:Nn \l__paper_code_tl {#1}
    \tl_replace_all:Nnn \l__paper_code_tl { / } { / \allowbreak }
    \tl_replace_all:Nnn \l__paper_code_tl { , } { , \allowbreak }
    \tl_replace_all:Nnn \l__paper_code_tl { \_ } { \_ \allowbreak }
    \tl_use:N \l__paper_code_tl
    \group_end:
  }
\ExplSyntaxOff

% ---------------------------------------------------------------- references
\usepackage[authoryear,round]{natbib}
\bibliographystyle{plainnat-amp} % plainnat with "&" in two-author in-text citations
\setlength{\bibsep}{2pt}

\usepackage[hidelinks,unicode]{hyperref}
\newcommand{\doi}[1]{doi:\href{https://doi.org/#1}{#1}}

% Running title (used in the page header of both documents)
\newcommand{\runningtitle}{PRISM: auditable, phase-aware glioma surveillance}
\newcommand{\papertitle}{PRISM: an auditable, phase-aware pipeline for post-treatment glioma surveillance ---
  rigorous evaluation of imaging, clinical and treatment-course information}
\fancyhead[L]{\small\itshape\runningtitle}
 directly after \documentclass[11pt,a4paper]{article}.
% Built with tectonic (XeTeX engine); also compiles with pdflatex.

% ---------------------------------------------------------------- fonts
\usepackage{iftex}
\ifPDFTeX
  \usepackage[T1]{fontenc}
  \usepackage{lmodern}
\else
  % XeTeX/LuaTeX: fontspec's default text font is Latin Modern (OpenType lmodern)
  \usepackage{fontspec}
\fi
\usepackage{textcomp}
\usepackage{microtype}
\usepackage[raggedright]{titlesec}

% ---------------------------------------------------------------- page layout
\usepackage[a4paper,margin=2.5cm,headheight=14pt]{geometry}
\linespread{1.3}
\usepackage{fancyhdr}
\pagestyle{fancy}
\fancyhf{}
\renewcommand{\headrulewidth}{0pt}
\fancyfoot[C]{\thepage}
\setlength{\parindent}{0pt}
\setlength{\parskip}{0.6\baselineskip}
\setlength{\emergencystretch}{3em}
\usepackage[running]{lineno} % continuous line numbers for reviewer comments
\renewcommand{\linenumberfont}{\normalfont\scriptsize\sffamily\color{gray}}

% ---------------------------------------------------------------- tables, figures
\usepackage{graphicx}
\graphicspath{{../docs/manuscript/figures/}} % single source: figures are not copied
\usepackage{booktabs}
\usepackage{array}
\usepackage{tabularx}
\usepackage{threeparttable}
\usepackage{pdflscape}
\usepackage{float}
\usepackage{placeins}
\usepackage{needspace}
\usepackage[titles]{tocloft}
\setlength{\cftsecnumwidth}{2.6em}
\setlength{\cftsubsecindent}{2.6em}
\setlength{\cftsubsecnumwidth}{3.4em}
\setlength{\cftbeforesecskip}{2pt}
\usepackage[font=small,labelfont=bf,labelsep=period,justification=justified,singlelinecheck=false]{caption}
\usepackage{siunitx}
\newcolumntype{L}{>{\raggedright\arraybackslash}X}
\newcolumntype{R}[1]{>{\raggedright\arraybackslash}p{#1}}
% U: like R, but \_ may break (plain-text variable names)
\newcolumntype{U}[1]{>{\raggedright\arraybackslash\renewcommand{\_}{\textunderscore\allowbreak}}p{#1}}
\setlength{\tabcolsep}{4pt}
\renewcommand{\arraystretch}{1.1}
\usepackage{etoolbox}
% floats (tables, figure legends) are single-spaced; running text keeps \linespread{1.3}
\AtBeginEnvironment{table}{\linespread{1}\selectfont}
\AtBeginEnvironment{figure}{\linespread{1}\selectfont}

% ---------------------------------------------------------------- colour, placeholders
\usepackage{xcolor}
% Visible placeholders: [CHECK ...] citations and [PENDING ...] author decisions.
\newcommand{\flag}[1]{\textcolor{red}{#1}}
\newcommand{\todo}[1]{\textcolor{red}{[#1]}}

% ---------------------------------------------------------------- inline code / file paths
% \code{...}: typewriter text that may break after "/", "_" and "," (robust, safe in captions).
% Inside \code, escape _ % & # { } as usual (\_ etc.).
\ExplSyntaxOn
\tl_new:N \l__paper_code_tl
\NewDocumentCommand{\code}{m}
  {
    \group_begin:
    \ttfamily\small
    \tl_set:Nn \l__paper_code_tl {#1}
    \tl_replace_all:Nnn \l__paper_code_tl { / } { / \allowbreak }
    \tl_replace_all:Nnn \l__paper_code_tl { , } { , \allowbreak }
    \tl_replace_all:Nnn \l__paper_code_tl { \_ } { \_ \allowbreak }
    \tl_use:N \l__paper_code_tl
    \group_end:
  }
\ExplSyntaxOff

% ---------------------------------------------------------------- references
\usepackage[authoryear,round]{natbib}
\bibliographystyle{plainnat-amp} % plainnat with "&" in two-author in-text citations
\setlength{\bibsep}{2pt}

\usepackage[hidelinks,unicode]{hyperref}
\newcommand{\doi}[1]{doi:\href{https://doi.org/#1}{#1}}

% Running title (used in the page header of both documents)
\newcommand{\runningtitle}{PRISM: auditable, phase-aware glioma surveillance}
\newcommand{\papertitle}{PRISM: an auditable, phase-aware pipeline for post-treatment glioma surveillance ---
  rigorous evaluation of imaging, clinical and treatment-course information}
\fancyhead[L]{\small\itshape\runningtitle}
 directly after \documentclass[11pt,a4paper]{article}.
% Built with tectonic (XeTeX engine); also compiles with pdflatex.

% ---------------------------------------------------------------- fonts
\usepackage{iftex}
\ifPDFTeX
  \usepackage[T1]{fontenc}
  \usepackage{lmodern}
\else
  % XeTeX/LuaTeX: fontspec's default text font is Latin Modern (OpenType lmodern)
  \usepackage{fontspec}
\fi
\usepackage{textcomp}
\usepackage{microtype}
\usepackage[raggedright]{titlesec}

% ---------------------------------------------------------------- page layout
\usepackage[a4paper,margin=2.5cm,headheight=14pt]{geometry}
\linespread{1.3}
\usepackage{fancyhdr}
\pagestyle{fancy}
\fancyhf{}
\renewcommand{\headrulewidth}{0pt}
\fancyfoot[C]{\thepage}
\setlength{\parindent}{0pt}
\setlength{\parskip}{0.6\baselineskip}
\setlength{\emergencystretch}{3em}
\usepackage[running]{lineno} % continuous line numbers for reviewer comments
\renewcommand{\linenumberfont}{\normalfont\scriptsize\sffamily\color{gray}}

% ---------------------------------------------------------------- tables, figures
\usepackage{graphicx}
\graphicspath{{../docs/manuscript/figures/}} % single source: figures are not copied
\usepackage{booktabs}
\usepackage{array}
\usepackage{tabularx}
\usepackage{threeparttable}
\usepackage{pdflscape}
\usepackage{float}
\usepackage{placeins}
\usepackage{needspace}
\usepackage[titles]{tocloft}
\setlength{\cftsecnumwidth}{2.6em}
\setlength{\cftsubsecindent}{2.6em}
\setlength{\cftsubsecnumwidth}{3.4em}
\setlength{\cftbeforesecskip}{2pt}
\usepackage[font=small,labelfont=bf,labelsep=period,justification=justified,singlelinecheck=false]{caption}
\usepackage{siunitx}
\newcolumntype{L}{>{\raggedright\arraybackslash}X}
\newcolumntype{R}[1]{>{\raggedright\arraybackslash}p{#1}}
% U: like R, but \_ may break (plain-text variable names)
\newcolumntype{U}[1]{>{\raggedright\arraybackslash\renewcommand{\_}{\textunderscore\allowbreak}}p{#1}}
\setlength{\tabcolsep}{4pt}
\renewcommand{\arraystretch}{1.1}
\usepackage{etoolbox}
% floats (tables, figure legends) are single-spaced; running text keeps \linespread{1.3}
\AtBeginEnvironment{table}{\linespread{1}\selectfont}
\AtBeginEnvironment{figure}{\linespread{1}\selectfont}

% ---------------------------------------------------------------- colour, placeholders
\usepackage{xcolor}
% Visible placeholders: [CHECK ...] citations and [PENDING ...] author decisions.
\newcommand{\flag}[1]{\textcolor{red}{#1}}
\newcommand{\todo}[1]{\textcolor{red}{[#1]}}

% ---------------------------------------------------------------- inline code / file paths
% \code{...}: typewriter text that may break after "/", "_" and "," (robust, safe in captions).
% Inside \code, escape _ % & # { } as usual (\_ etc.).
\ExplSyntaxOn
\tl_new:N \l__paper_code_tl
\NewDocumentCommand{\code}{m}
  {
    \group_begin:
    \ttfamily\small
    \tl_set:Nn \l__paper_code_tl {#1}
    \tl_replace_all:Nnn \l__paper_code_tl { / } { / \allowbreak }
    \tl_replace_all:Nnn \l__paper_code_tl { , } { , \allowbreak }
    \tl_replace_all:Nnn \l__paper_code_tl { \_ } { \_ \allowbreak }
    \tl_use:N \l__paper_code_tl
    \group_end:
  }
\ExplSyntaxOff

% ---------------------------------------------------------------- references
\usepackage[authoryear,round]{natbib}
\bibliographystyle{plainnat-amp} % plainnat with "&" in two-author in-text citations
\setlength{\bibsep}{2pt}

\usepackage[hidelinks,unicode]{hyperref}
\newcommand{\doi}[1]{doi:\href{https://doi.org/#1}{#1}}

% Running title (used in the page header of both documents)
\newcommand{\runningtitle}{PRISM: auditable, phase-aware glioma surveillance}
\newcommand{\papertitle}{PRISM: an auditable, phase-aware pipeline for post-treatment glioma surveillance ---
  rigorous evaluation of imaging, clinical and treatment-course information}
\fancyhead[L]{\small\itshape\runningtitle}

\renewcommand{\thesection}{S\arabic{section}}
\renewcommand{\thesubsection}{\thesection.\arabic{subsection}}
\captionsetup[table]{name=Supplementary Table}
\captionsetup[figure]{name=Supplementary Figure}
\fancyhead[R]{\small Supplementary Material}

\begin{document}
\linenumbers

\begin{flushleft}
{\Large\bfseries Supplementary Material\par}
\vspace{0.8em}
{\large\bfseries \papertitle\par}
\end{flushleft}

\emph{Draft, 2026-10-15. Every number below is copied from the result documents in
\code{docs/results/} and
\code{docs/manuscript/} or from
\code{outputs/results/}; the file is named in each section.
``Grade 4'' (G4)
is the primary population; ``all'' = all gliomas. AUC = IPCW AUC(120 d), mean over 5 repeats of the per-repeat
pooled out-of-fold value, with 95\% patient-cluster bootstrap CI (1,000 draws), unless stated otherwise.
PRISM v1 is the first version of PRISM, described in an unpublished internal report; ``PRISM's initial radiomics
model'' and ``our initial evaluation split'' are its radiomics model and its seed-62 test split. ``Baseline'' =
PRISM's core model (clinical + timing + as-of-scan treatment, ridge), the deployed configuration and the reference
for every imaging comparison.}

\section*{Contents}
\makeatletter\@starttoc{toc}\makeatother

\FloatBarrier
\section{Pre-specification register}\label{sec:S1}

\begin{table}[H]
\centering
\footnotesize
\begin{tabularx}{\linewidth}{@{}R{3.1cm} L R{3.0cm} R{3.0cm} R{3.3cm}@{}}
\toprule
Analysis & Status & Plan record & Decision rule / primary quantity & Outcome \\
\midrule
Radiomics vs baseline, nested over 70 configurations, grade 4 & Pre-specified primary (development) & Publication plan \S{}3.3 & $\Delta$AUC with 95\% CI & $-$0.032 [$-$0.084, +0.017]; no gain \\
LUMIERE locked clinical + timing + volumes, grade-4-trained & Pre-specified external model & \code{EXTERNAL\_LUMIERE.md} & Pooled IPCW AUC & 0.608 / 0.620 (ridge / boosting) \\
Phase-of-care ladder, post-RT R2 + I2 vs R2 & Pre-specified primary; plan in module docstring, not separately committed; a smoke run printed some I0--I2 numbers before the full run; plan unchanged & \code{PHASE\_LADDER.md} \S{}1 & $\Delta$AUC with CI and 100-permutation p & +0.061 [$-$0.038, +0.156], p = 0.21; not supported \\
Imaging-defined (clean) label & Pre-specified; plan snapshot saved before the run & \code{outputs/results/analyses/clean\_label/PLAN\_snapshot.md} & $\Delta$AUC $\geq$ 0.03 with CI excluding 0 & +0.027 [$-$0.033, +0.082]; not met \\
Re-segmentation & Rule written before any result was seen & \code{RESEGMENTATION.md} & Detection $\Delta$ $\geq$ 0.05 with CI > 0, or ET larger at progression in > 65\% of pairs & $-$0.016 [$-$0.057, +0.025]; 45\%; not met \\
CFB-GBM ADC & Pre-specified; plan snapshot saved before the run & \code{outputs/results/analyses/cfb\_adc/PLAN\_snapshot.md} & $\Delta$(d $-$ c) $\geq$ 0.03 with CI excluding 0 & +0.007 [$-$0.019, +0.035], p = 0.151; not met \\
Audit layer P1 (coverage) and P2 (abstention beyond case mix) & Pre-specified; plan snapshot saved before the run & \code{outputs/results/analyses/audit\_system/PLAN\_snapshot.md} & P1 coverage $\geq$ 0.78; P2 CI above 0 & P1 0.841 [0.790, 0.885], met; P2 +0.006 [$-$0.023, +0.036], not met \\
Split variance (500 splits) & Descriptive; designed after the single-split result was known & \code{SPLIT\_VARIANCE.md} & --- & --- \\
Positive controls, SHAP, feature conflicts, imbalance, schedule coupling, baselines, sensitivity, utility & Secondary or exploratory & respective documents & --- & --- \\
\bottomrule
\end{tabularx}
\end{table}

Sequence: the clean-label, re-segmentation, CFB-GBM and audit analyses were designed after the main imaging
result (no added value) was known. Each had its own rule fixed before its own run. Contrasts beyond the
primary ones are not corrected for multiplicity.

\FloatBarrier
\section{Cohort flow}\label{sec:S2}

Source: \code{docs/manuscript/FLOW.md}
(\code{outputs/data/mu\_flow.csv},
\code{outputs/results/analyses/tables/lumiere\_flow.csv}).
Figure: \code{docs/manuscript/cohort\_flow.png}
(Figure S1).

\begin{table}[H]
\centering
\small
\renewcommand{\thetable}{S1a}\renewcommand{\theHtable}{S1a}
\caption{\textbf{MU-Glioma-Post.}}
\label{tab:S1a}
\begin{tabularx}{\linewidth}{@{}L R{0.88cm} R{1.40cm} R{1.40cm}@{}}
\toprule
Step & Scans & Patients & Excluded \\
\midrule
All timepoints & 596 & 203 & -- \\
All four sequences and mask present & 594 & 203 & 2 \\
Tumour (labels 1--3) present & 592 & 203 & 2 \\
MRI day known & 589 & 202 & 3 \\
MRI day plausible ($\geq$ $-$365 d, not before an earlier timepoint) & 588 & 202 & 1 \\
Progression date known for progressors & 588 & 202 & 0 \\
MRI not after recorded death & 580 & 202 & 8 \\
Before late treatment (2nd-line, immunotherapy, new treatment) & 501 & 199 & 79 \\
\textbf{Before first progression (index scans)} & \textbf{308} & \textbf{155} & 193 \\
Grade 4 subset & 249 & 126 & -- \\
\bottomrule
\end{tabularx}
\end{table}

\begin{table}[H]
\centering
\small
\renewcommand{\thetable}{S1b}\renewcommand{\theHtable}{S1b}
\caption{\textbf{LUMIERE.}}
\label{tab:S1b}
\begin{tabularx}{\linewidth}{@{}L R{1.22cm} R{1.40cm} R{1.40cm}@{}}
\toprule
Step & Studies & Patients & Excluded \\
\midrule
All rated studies & 616 & 91 & -- \\
Post-operative (week $\geq$ 2 after surgery) & 456 & 81 & 160 \\
Rated SD / PR / CR / Post-Op & 199 & 70 & 257 \\
Before first progression & 132 & 61 & 67 \\
\textbf{Not after recorded death (index scans)} & \textbf{132} & \textbf{61} & 0 \\
\bottomrule
\end{tabularx}
\end{table}

Censored scans (52 MU, 39 grade 4, 9 LUMIERE at 120 d) stay in the cohort and are handled by the
censoring-aware labels and IPCW.

\textbf{Table S2.} The full cohort table with definitions is
\code{docs/manuscript/TABLE1\_COHORT.md}
(\code{table1.csv}). Additional rows not in main Table 1: last documented follow-up after scan, median [IQR]
days 275 [107--471] (all), 285 [112--478] (G4), 462 [238--721] (LUMIERE); time to progression or censoring 125
[50--237], 111 [48--208], 168 [91--287]; progression documented after the scan at any time 220 (71\%), 190
(76\%), 122 (92\%); LUMIERE T1c field strength 1.5 T / 3 T / other or missing 81 (61\%) / 44 (33\%) / 7; MU
field strength not distributed. WHO grade in all gliomas: grade 1, 3 (2\%); grade 2, 23 (15\%); grade 3, 2
(1\%); ``3 vs 4'', 1 (1\%); grade 4, 126 (81\%).

\FloatBarrier
\section{Outcome, censoring rules and label formulations}\label{sec:S3}

Source: \code{docs/results/FIX\_LABELS.md},
\code{docs/results/C1\_endpoint.md},
\code{docs/results/C1\_full\_tables.md}.

\textbf{Estimand.} Documented first progression within H = 120 d of the index scan (90 / 180 d as
sensitivities). Death without documented progression censors at the death day. Label 0 requires progression
after H or progression-free follow-up $\geq$ H. Composite (progression or death) is a sensitivity analysis
only.

\begin{table}[H]
\centering
\small
\renewcommand{\thetable}{S3}\renewcommand{\theHtable}{S3}
\caption{\textbf{Rules adopted in the label audit.}}
\label{tab:S3}
\begin{tabularx}{\linewidth}{@{}R{0.7cm} L L@{}}
\toprule
\# & Problem found & Rule \\
\midrule
1 & Death without documented progression treated as a verified negative at every horizon & Death censors at the death day (primary); composite event as sensitivity \\
1b & MRIs dated after the recorded death (15/96 MU deaths) & If more than one dated record follows the death day, the death date is ignored (7 patients, 3 with cohort scans); if exactly one, that scan is excluded \\
2 & Progression flag 0 with a progression date (3 patients); two date fields disagree (29 patients) & Any patient with a progression date is a progressor; primary date \code{date\_of\_first\_progression}; reverse order as sensitivity \\
3 & One MRI day of $-$965; progressor who died before progression & MRI day $\geq$ $-$365 and not before an earlier timepoint; death-before-progression flagged \\
4 & Hazard formulation counted scans censored mid-interval as surviving the interval & Censored scans contribute only fully observed intervals \\
5 & Follow-up extended by non-imaging dates & Kept as primary; imaging-only follow-up as sensitivity \\
6 & Unknown RT end coded as ``not early post-RT'' & \code{early\_post\_rt} is 1 / 0 / missing (94 scans missing, 65 grade 4) \\
\bottomrule
\end{tabularx}
\end{table}

Effect: six scans left the cohort (314 $\rightarrow$ 308; grade 4 253 $\rightarrow$ 249). At 120 d, 17
negatives became censored (16 deaths within 120 d without progression, 1 ignored death date) and 1 negative
became positive; no positive changed.

\begin{table}[H]
\centering
\small
\renewcommand{\thetable}{S4}\renewcommand{\theHtable}{S4}
\caption{\textbf{Label counts (positive / verified negative / censored).}}
\label{tab:S4}
\begin{tabularx}{\linewidth}{@{}R{3.42cm} L R{2.38cm} R{2.55cm} R{2.55cm}@{}}
\toprule
Population & Label & H = 90 & H = 120 & H = 180 \\
\midrule
All (308 / 155) & primary & 65 / 195 / 48 & 100 / 156 / 52 & 144 / 105 / 59 \\
All & composite & 79 / 195 / 34 & 116 / 156 / 36 & 161 / 105 / 42 \\
All & alternative progression date & 63 / 191 / 54 & 99 / 151 / 58 & 141 / 102 / 65 \\
All & imaging-only follow-up & 65 / 187 / 56 & 100 / 149 / 59 & 144 / 99 / 65 \\
Grade 4 (249 / 126) & primary & 62 / 152 / 35 & 93 / 117 / 39 & 131 / 75 / 43 \\
Grade 4 & composite & 74 / 152 / 23 & 107 / 117 / 25 & 146 / 75 / 28 \\
LUMIERE (132 / 61) & primary & 24 / 101 / 7 & 47 / 76 / 9 & 62 / 61 / 9 \\
LUMIERE & composite & 26 / 101 / 5 & 51 / 76 / 5 & 66 / 61 / 5 \\
\bottomrule
\end{tabularx}
\end{table}

LUMIERE gives a survival time without a death indicator; a recorded survival time is read as a death and used
only as a censoring time in the primary label.

\textbf{IPCW AUC.} Cumulative/dynamic AUC at 120 d with Kaplan--Meier censoring weights
\citep{uno2007,blanche2013}. Cases are scans with an event by 120 d, weighted by the inverse of the censoring
survival at their event time; controls are scans event-free beyond 120 d. The vectorised implementation
matches scikit-survival to \num{2e-16} (\code{FIX\_EVALUATION.md} \S{}1). In the bootstrap, censoring
weights are re-estimated in each resample; if a resample has no censoring, G = 1.

\begin{table}[H]
\centering
\small
\renewcommand{\thetable}{S5}\renewcommand{\theHtable}{S5}
\caption{\textbf{Label formulations, grade 4, H = 120 d} (\code{C1\_endpoint.md}; $\Delta$AUC vs naive).}
\label{tab:S5}
\begin{tabularx}{\linewidth}{@{}L R{1.92cm} R{3.60cm} R{4.12cm} R{0.80cm} R{0.88cm}@{}}
\toprule
Features & Formulation & AUC [95\% CI] & $\Delta$AUC vs naive [95\% CI] & O/E & Slope \\
\midrule
volumes + clinical & naive & 0.651 [0.567, 0.727] & --- & 1.14 & 0.95 \\
 & verified & 0.657 [0.577, 0.732] & +0.006 [$-$0.020, +0.030] & 0.95 & 0.91 \\
 & IPCW & 0.656 [0.578, 0.731] & +0.005 [$-$0.021, +0.029] & 0.97 & 0.92 \\
 & hazard & 0.653 [0.574, 0.726] & +0.002 [$-$0.018, +0.022] & 0.87 & 0.50 \\
legacy radiomics + clinical & naive & 0.648 [0.553, 0.736] & --- & 1.13 & 0.52 \\
 & verified & 0.655 [0.575, 0.730] & +0.007 [$-$0.021, +0.038] & 0.96 & 0.51 \\
 & IPCW & 0.654 [0.575, 0.730] & +0.006 [$-$0.022, +0.036] & 0.98 & 0.51 \\
 & hazard & 0.669 [0.583, 0.746] & +0.021 [$-$0.027, +0.068] & 0.90 & 0.33 \\
\bottomrule
\end{tabularx}
\end{table}

All 36 $\Delta$AUC CIs (2 populations $\times$ 2 feature sets $\times$ 3 horizons $\times$ 3 formulations)
include 0. The naive label under-predicts risk in all 12 settings (O/E 1.12--1.19 against 0.95--0.99 for
verified / IPCW). The legacy radiomics pipeline in C1 predates the clipping fix (S5).

\FloatBarrier
\section{Predictor definitions and the as-of-scan treatment rule}\label{sec:S4}

Source: \code{prism/features/tabular.py};
\code{docs/results/FIX\_TREATMENT\_ASOF.md};
\code{docs/results/FIX\_FEATURES.md}; audit record input tables
(\code{outputs/results/analyses/audit\_system/records.md}).

\begin{table}[H]
\centering
\small
\renewcommand{\thetable}{S6}\renewcommand{\theHtable}{S6}
\caption{\textbf{Baseline predictors (23 raw; 31 design columns after missingness indicators).}}
\label{tab:S6}
\begin{tabularx}{\linewidth}{@{}R{1.6cm} U{3.4cm} L R{4.6cm}@{}}
\toprule
Group & Predictor & Definition (as of scan day) & Source column(s) \\
\midrule
Clinical & age & age at diagnosis, years & \code{clinical\_age\_at\_diagnosis} \\
 & male & sex at birth & \code{clinical\_sex\_at\_birth} \\
 & idh\_mutant, idh\_unknown & IDH1/IDH2 mutant; neither documented & \code{clinical\_idh1\_mutation}, \code{clinical\_idh2\_mutation} \\
 & mgmt\_methylated, mgmt\_unknown & MGMT promoter methylated; unknown code & \code{clinical\_mgmt\_methylation} \\
Timing & days\_from\_diagnosis & scan day & \code{days\_from\_diagnosis\_to\_mri} \\
 & days\_since\_rt\_end & days since documented RT end; missing until RT completed & RT end day vs scan day \\
 & rt\_date\_unknown & no completed RT with known end date at the scan & RT end day \\
 & early\_post\_rt & 0--84 d after RT end (RANO 12-week window); missing if RT end unknown & RT end day \\
 & early\_post\_rt\_x\_mgmt\_methylated & interaction & RT end day $\times$ MGMT \\
 & visit\_index & rank of the scan among the patient's MRIs & MRI days \\
Treatment & rt\_started & RT start day $\leq$ scan day & \code{clinical\_radiation\_therapy} + RT start day \\
 & rt\_completed & RT end day $\leq$ scan day & RT end day \\
 & rt\_dose\_gy, rt\_fractions, rt\_hypofractionated & prescription, visible once RT started; hypofractionation $\leq$ 15 fractions & \code{clinical\_dose}, \code{clinical\_number\_of\_fractions} \\
 & rt\_start\_day & days from diagnosis to RT start, visible once started & RT start day \\
 & rt\_duration\_days & RT end $-$ start, visible once completed & RT start and end days \\
 & chemo\_started & initial TMZ / lomustine start $\leq$ scan day & initial chemotherapy fields \\
 & surgery\_day, biopsy\_before\_resection & visible once the first surgery has happened & first surgery fields \\
 & previous\_brain\_tumor & history, known at any scan & \code{clinical\_previous\_brain\_tumor} \\
\bottomrule
\end{tabularx}
\end{table}

Unknown values are imputed in-fold with the training median plus a missingness indicator; nothing is
zero-filled silently. Therapy marked given without a start date (2 RT, 2 chemotherapy patients) is assumed to
have started before the scan.

\begin{table}[H]
\centering
\small
\renewcommand{\thetable}{S7}\renewcommand{\theHtable}{S7}
\caption{\textbf{What leaked before the as-of-scan fix} (scans affected, grade 4 / all gliomas).}
\label{tab:S7}
\begin{tabularx}{\linewidth}{@{}R{4.6cm} L R{4.2cm}@{}}
\toprule
Feature & Problem & Affected \\
\midrule
\code{rt\_given}, \code{chemo\_given} & ``will receive RT / chemotherapy'', true for scans before the course started & 13 / 23 (RT), 13 / 21 (chemo) \\
dose, fractions, hypofractionation, \code{rt\_start\_day} & prescription and start date of a course that had not started & 13 / 19--23 \\
\code{rt\_duration\_days} & duration of a course not yet finished (also reveals early stopping) & 17 / 29 \\
\code{days\_since\_rt\_end}, \code{rt\_date\_unknown} & negative values from the future RT end date & 17 / 29 \\
\code{surgery\_day}, \code{biopsy\_before\_resection} & one scan precedes the first surgery & 1 / 1 \\
\bottomrule
\end{tabularx}
\end{table}

Progression, death, last follow-up and late-treatment dates are used only for labels, censoring and
eligibility. Change features use earlier scans only.

\textbf{Effect on repeated CV (intermediate inputs; earlier repeat-averaged estimator;
\code{FIX\_TREATMENT\_ASOF.md} \S{}3).} Baseline ridge grade 4 0.694 [0.613, 0.775]
$\rightarrow$ 0.729
[0.648, 0.811]; all gliomas 0.719 [0.648, 0.788] $\rightarrow$ 0.749 [0.675, 0.820]. Boosting grade 4 0.723
[0.633, 0.806] $\rightarrow$ 0.731 [0.641, 0.813]; all gliomas 0.743 [0.666, 0.816] $\rightarrow$ 0.740
[0.663, 0.815]. Timing-only boosting fell (grade 4 0.718 $\rightarrow$ 0.696; all 0.735 $\rightarrow$ 0.705):
part of its earlier signal came from the future RT end date. The MGMT $\times$ window flag changed repeated-CV
AUC by $\leq$ 0.008. These intermediate values are superseded by the final run (main Table 2).

\textbf{Molecular look-ahead found by the audit (exploratory).} IDH and MGMT are used at scans taken before
the diagnosis date: 8 grade 4 scans (7 patients; 1 event, 4 non-events, 3 censored) and 11 all-glioma scans,
1--58 days before diagnosis (\code{AUDIT\_SYSTEM.md} \S{}2). Its effect on the baseline AUC was not
quantified.

\textbf{Imaging measurements.} Volumes (ET, NETC, SNFH, RC, tumour = ET + NETC, lesion), morphology (ET
fraction, multifocality, largest axial ET area, ET-to-cavity distance and proximity, T1c / FLAIR intensities
relative to brain > 10 mm (Euclidean) from lesion and cavity), and past-only change (from previous scan and
nadir, growth rate, \code{vol\_pd\_flag} requiring a measurable nadir $\geq$ 0.1 cc,
\code{et\_new\_from\_zero\_nadir}, missingness flags for
undefined relative change; 25 change columns). For
LUMIERE transport, necrosis and cavity were merged
(\code{log\_necrosis\_cavity\_cc}) because DeepBraTumIA
labels the cavity as necrosis (\code{FIX\_FEATURES.md} \S{}1).

\FloatBarrier
\section{Radiomics: extraction, configurations and outlier clipping}\label{sec:S5}

Source: \code{docs/results/C3\_radiomics.md},
\code{docs/results/FIX\_RADIOMICS\_CLIP.md},
\code{docs/results/FIX\_FEATURES.md} \S{}6,
\code{prism/features/radiomics.py}.

\textbf{Extraction.} PyRadiomics \citep{vangriethuysen2017}. \emph{Legacy} features are those of PRISM's
initial radiomics model: T1c and FLAIR, original, LoG and wavelet images, 2,446 columns. Under the legacy within-lesion z-score
and fixed bin width the texture matrices are effectively two-level (GLCM joint entropy $\leq$ 2 bits), so
legacy texture behaves as a binary habitat descriptor (above vs below lesion mean). \emph{Corrected} (v2)
features use a whole-brain z-score, a fixed bin count of 32 (joint entropy 4.7--6.3 bits), no second N4
correction, and are extracted per compartment (ET, NETC, SNFH, core, lesion, 5 mm peri-cavity rim, shape) and
for all four sequences (2,316 columns). Known limitations: a fixed bin count rather than the IBSI-recommended
fixed bin width \citep{zwanenburg2020}; the ICC stability pass covers T1c and FLAIR texture and first-order
features only (1,116 of 2,316 columns); the 5 mm rim uses a city-block dilation.

\textbf{Configurations (70).} Variant (legacy; corrected whole lesion; corrected per compartment; corrected
all four sequences) $\times$ stability filter (none; ICC $\geq$ 0.75 under 1-voxel erosion, in-fold) $\times$
selection (top-10 or top-30 univariate F; L1 with non-zero coefficients only, at most 30; PCA-10;
decorrelation + top-20) $\times$ integration (concatenation with the baseline; or a stacked radiomics score
from inner out-of-fold predictions). Nested choice: inner 3-fold patient-grouped AUC in each outer training
fold. Missing compartments are imputed with the median and one indicator per missingness pattern.

\textbf{Outlier clipping.} Three legacy columns (FLAIR LoG $\sigma$ = 1, 2, 3 mm NGTDM Coarseness) held the
PyRadiomics sentinel $10^{6}$ in four scans (two grade 4); other values in these columns have medians of
$10^{-4}$ to \num{8e-4}. One such value inflates a training SD by orders of magnitude and makes PCA and the
stacked score depend on these scans. Before the fix, the stacked score had inner-OOF training values of
$-10^{4}$ to $-$\num{4.6e6} in 20/25 grade 4 fits. The fix, \code{QuantileClip(q = 0.01)}, is fitted in each
training (and inner) fold, clips to the fold's 1st/99th order statistics, and applies the same bounds to test
scans. It does not look at labels and applies to all configurations and to deep features.

\begin{table}[H]
\centering
\footnotesize
\renewcommand{\thetable}{S8}\renewcommand{\theHtable}{S8}
\caption{\textbf{Effect of the clipping fix} (\code{FIX\_RADIOMICS\_CLIP.md}).}
\label{tab:S8}
\begin{tabularx}{\linewidth}{@{}R{1.9cm} L L R{1.7cm}@{}}
\toprule
Population & Nested AUC before $\rightarrow$ after & $\Delta$AUC before $\rightarrow$ after & Cal. slope \\
\midrule
Grade 4 & 0.672 [0.597, 0.743] $\rightarrow$ 0.666 [0.586, 0.739] & $-$0.026 [$-$0.073, +0.017] $\rightarrow$ $-$0.032 [$-$0.084, +0.017] & 0.50 $\rightarrow$ 0.60 \\
All gliomas & 0.732 [0.664, 0.794] $\rightarrow$ 0.747 [0.678, 0.811] & +0.012 [$-$0.029, +0.050] $\rightarrow$ +0.028 [$-$0.017, +0.071] & 0.77 $\rightarrow$ 0.85 \\
\bottomrule
\end{tabularx}
\end{table}

Largest change in any single configuration: 0.020. Single configurations with CI above 0: grade 4 0/70 before
and after (below 0: 12/70 $\rightarrow$ 8/70); all gliomas 0/70 $\rightarrow$ 1/70
(\code{legacy|none|pca10|stack|ridge}, 0.754, +0.035 [+0.002, +0.068], p($\Delta$ $\leq$ 0) = 0.021). Not
covered by the fix: the radiomics group of the feature-group search (S10), the C1 radiomics feature set (S3),
and the feature-conflict B + radiomics row (S12).

\begin{landscape}
\begin{table}[H]
\centering
\small
\renewcommand{\thetable}{S9}\renewcommand{\theHtable}{S9}
\caption{\textbf{Radiomics results by population} (\code{C3\_radiomics.md};
\code{outputs/results/c3/\{grade4,all\_glioma\}/results.csv}).}
\label{tab:S9}
\begin{tabularx}{\linewidth}{@{}L R{3.6cm} R{4.3cm} R{1.0cm} R{3.6cm} R{4.3cm} R{1.0cm}@{}}
\toprule
Model & Grade 4 AUC & $\Delta$ [95\% CI] & Slope & All gliomas AUC & $\Delta$ [95\% CI] & Slope \\
\midrule
Baseline ridge & 0.698 [0.615, 0.774] & --- & 0.89 & 0.720 [0.638, 0.792] & --- & 0.88 \\
Baseline boosting & 0.687 [0.601, 0.768] & $-$0.012 [$-$0.072, +0.048] & 0.82 & 0.718 [0.642, 0.792] & $-$0.002 [$-$0.052, +0.047] & 0.87 \\
Best post hoc (legacy PCA-10 stack) & 0.721 [0.638, 0.793] & +0.023 [$-$0.013, +0.057] & 0.92 & 0.754 [0.685, 0.823] & +0.035 [+0.002, +0.068] & 0.96 \\
\textbf{Nested over 70 (honest)} & \textbf{0.666 [0.586, 0.739]} & \textbf{$-$0.032 [$-$0.084, +0.017]} & 0.60 & \textbf{0.747 [0.678, 0.811]} & \textbf{+0.028 [$-$0.017, +0.071]} & 0.85 \\
\bottomrule
\end{tabularx}
\end{table}
\end{landscape}

Median $\Delta$AUC by variant, grade 4: legacy +0.005, v2 all four sequences $-$0.016, v2 compartments
$-$0.017, v2 lesion $-$0.001; all gliomas +0.024, +0.000, $-$0.005, +0.005. Nested choices: legacy in 18/25
(grade 4) and 25/25 (all) outer folds. Stacking beat concatenation in grade 4 (median $\Delta$ $-$0.003 vs
$-$0.027).

\FloatBarrier
\section{Frozen deep features}\label{sec:S6}

Source: \code{docs/results/DEEP\_FEATURES.md};
\code{outputs/results/analyses/deep/}.

\textbf{Encoders.} MONAI SwinUNETR (four channels FLAIR / T1c / T1 / T2, feature size 48), encoder only: BraTS
2021 weights; BrainSegFounder (UK Biobank self-supervised pretraining, BraTS fine-tuning; \citealp{cox2024}
\flag{[CHECK]}); and random weights (seed 0) as the control. \textbf{Input.} Skull-stripped 1 mm volumes,
z-scored within the brain, reoriented to LPS, zero-padded to a multiple of 32; eight flip test-time
augmentations. \textbf{Features.} Mask-weighted mean pooling at stages 0--3 (48 + 96 + 192 + 384 = 720 per
region); MU regions lesion, ET, NETC, SNFH, 5 mm rim and ET $\cup$ SNFH; LUMIERE regions ET, SNFH, ET $\cup$
SNFH and lesion. The 768-d bottleneck was not used (between-scan SD \num{3e-4} for BraTS 2021 weights).
\textbf{Families.} D1 pretrained, 32 configurations (2 encoders $\times$ \{lesion, compartments\} $\times$ 4
selections $\times$ 2 integrations); D0 random, 16; D1b deep + legacy radiomics stacked, 16; DS shared
regions, 16. \textbf{Licences.} BraTS 2021 SwinUNETR weights Apache-2.0; BrainSegFounder weights under the UK
Biobank material transfer agreement (not redistributed).

\begin{table}[H]
\centering
\footnotesize
\renewcommand{\thetable}{S10}\renewcommand{\theHtable}{S10}
\caption{\textbf{Deep features on MU.}}
\label{tab:S10}
\begin{tabularx}{\linewidth}{@{}L R{2.75cm} R{3.5cm} R{2.75cm} R{3.5cm}@{}}
\toprule
Arm & Grade 4 AUC & Grade 4 $\Delta$ [95\% CI] & All AUC & All $\Delta$ [95\% CI] \\
\midrule
Baseline & 0.698 [0.615, 0.774] & --- & 0.720 [0.638, 0.792] & --- \\
D1 pretrained, nested over 32 & 0.694 & $-$0.004 [$-$0.069, +0.063] & 0.731 & +0.011 [$-$0.032, +0.056] \\
D0 random weights, nested over 16 & 0.649 & $-$0.049 [$-$0.085, $-$0.014] & 0.711 & $-$0.009 [$-$0.032, +0.009] \\
D1b deep + legacy radiomics, nested over 16 & 0.709 & +0.011 [$-$0.032, +0.054] & 0.751 & +0.031 [$-$0.003, +0.066] \\
DS shared regions, nested over 16 & 0.699 & +0.001 [$-$0.057, +0.060] & 0.725 & +0.005 [$-$0.035, +0.046] \\
Baseline + volumes + D1 vs baseline + volumes & 0.694 vs 0.702 & $-$0.008 [$-$0.045, +0.033] & 0.739 vs 0.733 & +0.007 [$-$0.021, +0.036] \\
\bottomrule
\end{tabularx}
\end{table}

Single configurations with CI above 0: 0/80 (grade 4); 3/80 (all gliomas; +0.035 to +0.040, lower bounds
+0.001 to +0.004); about 2 of 80 expected by chance. Image-only stacked score, grade 4: pretrained 0.652
(modal configuration), random 0.566; correlation with log tumour volume r = 0.58 and with the baseline score r
= 0.28. Calibration slope of nested D1: 0.55 (grade 4), 0.79 (all).

\FloatBarrier
\section{PRISM v1: the initial radiomics model, the initial evaluation split and the fusion components}\label{sec:S7}

Source: \code{docs/results/SHAP\_PRISM\_SPLIT.md};
\code{docs/manuscript/FIGURES.md} (Figure 1);
\code{outputs/results/analyses/
shap\_prism\_split/};
\code{outputs/results/analyses/paper\_figures/single\_split\{,\_pairs\}.csv}.

\textbf{Reconstruction.} The input matrix was rebuilt from the legacy feature files plus the clinical codes
using PRISM v1's case table; the shipped bundle reproduced the stored test predictions on all 84 rows
(maximum absolute difference \num{1.1e-15} calibrated, \num{1.4e-15} raw). The model is an L2 logistic
regression (liblinear, C = 0.0087, balanced class weights) on 46 standardised columns (two of 48 selected
columns are dropped by a correlation filter), with a monotone Platt calibrator. Exact linear SHAP reproduced
the shipped importance table (Spearman 0.997, identical top 10). Modality shares: T1c 52\%, FLAIR 26\%,
clinical codes 22\% (share of mean absolute per-scan group sum).

\begin{table}[H]
\centering
\small
\renewcommand{\thetable}{S11}\renewcommand{\theHtable}{S11}
\caption{\textbf{Composition of the test set of our initial evaluation split} (percentile among 500 design splits).}
\label{tab:S11}
\begin{tabularx}{\linewidth}{@{}L R{2.02cm} R{2.55cm} R{3.42cm} R{1.75cm}@{}}
\toprule
 & Full cohort & Initial split test set & 500 splits: mean (2.5--97.5\%) & Percentile \\
\midrule
Scans / patients & 308 / 155 & 82 / 29 & 84.5 / 29.8 &  \\
Event rate, naive & 0.325 & 0.268 & 0.254 (0.165--0.348) & 62 \\
Grade 4 / GBM & 0.81 / 0.70 & 0.80 / 0.77 & 0.82 / 0.63 & 44 / 97 \\
IDH-mutant scans & 0.195 & 0.146 & 0.271 (0.131--0.403) & 5 \\
Pre-RT scans & 0.373 & 0.439 & 0.353 (0.220--0.474) & 89 \\
Event rate in pre-RT scans & 0.087 & 0.000 (0 / 36) & 0.044 (0--0.138) & 15 \\
$\leq$ 12 weeks post-RT & 0.338 & 0.244 & 0.283 & 16 \\
Later post-RT & 0.263 & 0.268 & 0.345 (0.243--0.450) & 6 \\
AUC of the rule ``scan is post-RT'' & 0.677 & 0.740 & 0.675 (0.575--0.772) & 92 \\
\bottomrule
\end{tabularx}
\end{table}

\begin{table}[H]
\centering
\small
\renewcommand{\thetable}{S12}\renewcommand{\theHtable}{S12}
\caption{\textbf{Pair decomposition (naive AUC)} (\code{single\_split\_pairs.csv}; 22
positives, 36 pre-RT negatives, 24 other negatives; 60\% of pairs involve a pre-RT negative).}
\label{tab:S12}
\begin{tabularx}{\linewidth}{@{}L R{1.57cm} R{2.20cm} R{2.20cm}@{}}
\toprule
Model & All pairs & Positives vs pre-RT negatives & Positives vs other negatives \\
\midrule
PRISM initial radiomics model & 0.802 & 0.823 & 0.771 \\
Baseline ridge & 0.871 & 0.962 & 0.735 \\
Baseline GBM & 0.844 & 0.933 & 0.710 \\
Baseline + legacy radiomics & 0.862 & 0.934 & 0.754 \\
Clinical only (ridge) & 0.791 & 0.789 & 0.793 \\
\bottomrule
\end{tabularx}
\end{table}

On the 36 pre-RT scans the initial radiomics model's T1c features lowered the log-odds by $-$0.50 on average; the
baseline ridge did the same through treatment ($-$0.38) and timing ($-$0.15). Our models' SHAP profiles on
this split match their whole-cohort out-of-fold profiles at group level (baseline ridge: treatment 55\%,
clinical 29\%, timing 16\% on the split; 52 / 29 / 19\% out of fold); the split does not lean on unusual
features.

\textbf{Post-RT test scans} (42 scans, 19 patients, 19 positives; IPCW AUC [95\% CI], naive in brackets):
initial radiomics model 0.770 [0.604, 0.907] (0.778); baseline ridge 0.718 [0.481, 0.880] (0.730); baseline GBM
0.699 [0.463, 0.846] (0.705); + legacy radiomics 0.758 [0.605, 0.886] (0.769); clinical only 0.818 [0.664,
0.924] (0.826), $\Delta$ vs baseline +0.100 [$-$0.015, +0.293].

\textbf{Before and after the as-of-scan fix} (30 test patients / 84 scans, naive / IPCW;
\code{head\_to\_head\_\{before,after\}.csv}): initial radiomics
model 0.802 / 0.811 in both; baseline ridge 0.870 /
0.844 $\rightarrow$ 0.880 / 0.860; baseline boosting 0.901 / 0.881 $\rightarrow$ 0.869 / 0.855; baseline +
legacy radiomics 0.869 / 0.848 $\rightarrow$ 0.871 / 0.852; baseline + corrected compartments 0.869 / 0.843
$\rightarrow$ 0.886 / 0.870. With the final labels (29 / 82): ridge 0.871 / 0.851, boosting 0.844 / 0.824.

\textbf{PRISM v1 encoder and fusion components} (source: \code{legacy/reproduction/\{README,FINDINGS,AUDIT\}.md};
v1's own case table, labels and scoring on the 30-patient / 84-scan test set). All v1 numbers reproduced from the
code. The v1 image encoder was a small SwinUNETR-style transformer with random, untrained weights; its embeddings
alone gave AUC about 0.50 in fusion (linear probe 0.55 in training CV, 0.62 on the test set). The v1 late-fusion MLP
scored AUC 0.8006, below its radiomics input alone (0.8043); its checkpoint had been selected on test loss and its
0.48 threshold hand-tuned, so its reported recall and F1 gains over the radiomics model were an operating-point
shift. A BraTS-pretrained SwinUNETR fine-tuned on MU (lesion-masked pooling) reached image-only test AUC 0.64--0.65,
and late fusion with the radiomics score gave 0.790--0.809, with a patient-bootstrap CI of the difference spanning 0.
These results motivated the frozen-encoder design with a random-weight control (S6) and the single regularised model
with ablatable branches (main \S2.1). The v1 report generator's RANO section was computed from the radiomics score,
not from mask measurements; the rebuilt generator is described in S19.

\FloatBarrier
\section{Split variance: full tables and verifier checks}\label{sec:S8}

Source: \code{docs/results/SPLIT\_VARIANCE.md};
\code{docs/manuscript/FIGURES.md} \S{}2 (clipped
\code{baseline\_rad});
\code{outputs/results/analyses/split\_variance/\{summary,summary\_clipped,per\_split\_auc\_clipped\}.csv}.

\textbf{Design.} Seeds 0--499 of PRISM v1's split search gave 500 distinct test sets (82--94 test
scans, mean 86.1; 10--33 positives). The search favours patients with many scans: all 157 patients of
PRISM v1's case table appear, and the most frequent is in 71\% of splits. Cohort, label and features are
PRISM v1's (315 scans / 157 patients; naive label, equal to ours on all 314 shared scans). Rank = share of
splits with strictly lower AUC than the model's own seed-62 value re-run here.

\begin{table}[H]
\centering
\footnotesize
\renewcommand{\thetable}{S13}\renewcommand{\theHtable}{S13}
\caption{\textbf{Naive AUC over 500 splits.}}
\label{tab:S13}
\begin{tabularx}{\linewidth}{@{}L R{2.1cm} R{1.7cm} R{1.6cm} R{1.1cm} R{1.4cm} R{1.5cm}@{}}
\toprule
Model & Mean $\pm$ SD & 2.5--97.5\% & Seed 62 (rank) & $\geq$ 0.802 & Beats baseline & Mean $\Delta$ vs baseline \\
\midrule
Initial radiomics pipeline, re-selected per split & 0.720 $\pm$ 0.055 & 0.610--0.819 & 0.801 (92\%) & 7.2\% & 12\% & $-$0.081 \\
Initial radiomics model, shipped features refitted (optimistic) & 0.815 $\pm$ 0.044 & 0.726--0.887 & 0.802 (38\%) & 62\% & 58\% & +0.014 \\
Baseline ridge & 0.801 $\pm$ 0.052 & 0.694--0.895 & 0.871 (91\%) & 50\% & -- & -- \\
Baseline + legacy radiomics (clipped) & 0.799 $\pm$ 0.053 & 0.674--0.892 & 0.862 (89\%) & 50.2\% & 51.6\% & $-$0.002 \\
\bottomrule
\end{tabularx}
\end{table}

\begin{table}[H]
\centering
\small
\renewcommand{\thetable}{S14}\renewcommand{\theHtable}{S14}
\caption{\textbf{IPCW AUC over 475 splits.}}
\label{tab:S14}
\begin{tabularx}{\linewidth}{@{}L R{4.82cm} R{2.55cm} R{2.20cm} R{1.15cm}@{}}
\toprule
Model & Mean $\pm$ SD & Seed 62 (rank) & Beats baseline & Mean $\Delta$ \\
\midrule
Initial radiomics pipeline, re-selected & 0.727 $\pm$ 0.058 (0.602--0.832) & 0.816 (94\%) & 25\% & $-$0.042 \\
Shipped features refitted (optimistic) & 0.825 $\pm$ 0.045 (0.737--0.901) & 0.817 (40\%) & 81\% & +0.055 \\
Baseline ridge & 0.769 $\pm$ 0.056 (0.654--0.876) & 0.851 (92\%) & -- & -- \\
Baseline + legacy radiomics (clipped) & 0.774 $\pm$ 0.058 & 0.845 (92.0\%) & 56.8\% & +0.005 \\
\bottomrule
\end{tabularx}
\end{table}

Per split, clipped minus pre-clip naive AUC of the radiomics arm: mean $-$0.002, SD 0.009. Gap between honest
and optimistic versions of the initial radiomics pipeline: 0.095 $\pm$ 0.040.

\textbf{Composition.} Baseline ridge per-split AUC correlates with the pre-RT event rate (Spearman $-$0.61
naive, $-$0.57 IPCW), the AUC of the ``post-RT'' rule (+0.57, +0.51), the later post-RT share ($-$0.26) and
the pre-RT share (+0.22). Nine composition variables explain CV R$^{2}$ 0.42 and predict 0.851 for seed 62
(observed 0.871, mean 0.801). For the initial radiomics pipeline: CV R$^{2}$ 0.10 (naive), 0.15 (IPCW); strongest
correlates IDH-mutant fraction ($-$0.30) and the ``post-RT'' rule (+0.22); predicted 0.773 for seed 62
(observed 0.801, mean 0.720).

\textbf{Verifier checks (earlier per-split table; not repeated on the final run).} (1) On seed 62 and 15 other
splits, redoing the ranking at full cost (10 folds $\times$ 100 permutation repeats) gave 0.718 $\pm$ 0.054 on
the 15 splits (cheap ranking 0.726 $\pm$ 0.057; cheap $-$ full +0.008 $\pm$ 0.033, r = 0.83). Re-running the
shipped Optuna search as well gave 0.654 $\pm$ 0.054 ($-$0.064 $\pm$ 0.055 against full ranking with the
shipped C); its inner-CV objective averaged 0.865 (optimism $\approx$ 0.21), and on seed 62 it gave 0.762, not
0.802. (2) Against 300 fully random 30-patient splits, the design splits had more test scans (85.5 vs 59.9)
and a narrower baseline AUC spread (SD 0.054 vs 0.069), so spread under the design underestimates what a fresh
split would show \citep{varoquaux2018}.

\FloatBarrier
\section{Estimator and evaluation details}\label{sec:S9}

Source: \code{docs/results/FIX\_EVALUATION.md}.

\begin{itemize}
\item \textbf{Headline estimator.} Mean over the 5 repeats of the per-repeat pooled-OOF IPCW AUC (and Brier,
O/E, slope). The earlier estimator --- the AUC of the OOF prediction averaged over repeats --- is an ensemble
of five models per scan and is optimistic, more so for unstable nested pipelines. On intermediate inputs, the
nested radiomics $\Delta$AUC moved from +0.009 (repeat-averaged) to $-$0.002 (mean over repeats) under the
change.
\item \textbf{Bootstrap.} Patient-cluster bootstrap, 1,000 draws; each draw resamples patients once and
recomputes the mean-over-repeats statistic of every model on that draw, so differences are paired. $\Delta$AUC
is the point estimate (not the bootstrap mean) with the bootstrap CI.
\item \textbf{Calibration.} O/E (Kaplan--Meier risk / mean prediction, or IPCW-weighted),
calibration-in-the-large with a logit offset, calibration slope by IPCW-weighted logistic regression on the
logit, flexible calibration curves (restricted cubic spline, 3 knots) and ICI \citep{vancalster2019}.
\item \textbf{External estimator.} Every LUMIERE arm is scored on the same 5 $\times$ 5 patient-grouped
LUMIERE folds (per-fold IPCW AUC averaged over 25 folds, censoring weights from the whole LUMIERE cohort);
locked models are also scored pooled. An earlier pooled score of differently recalibrated folds produced an
artefactual AUC of 0.414 for recalibrated ridge; recalibration leaves discrimination unchanged by
construction.
\item \textbf{Selection inside folds.} Feature selection, PCA, ICC filtering, stacking and hyperparameter
choice are all inside the training fold, with inner patient-grouped CV for every choice
\citep{varma2006,cawley2010,demircioglu2021}.
\end{itemize}

\FloatBarrier
\section{Feature-group search with permutation control}\label{sec:S10}

Source: \code{docs/results/C2\_feature\_search.md};
\code{docs/results/FINAL\_RESULTS.md}.

Six groups --- clinical, timing, volumes, morphology, change, legacy radiomics (top 20 in-fold; unclipped) ---
give 63 combinations $\times$ 2 learners (ridge; fixed shallow boosting). The treatment block is not part of
this search. The honest estimate selects one of the 126 candidates per outer fold by inner CV. The permutation
control moves each patient's whole outcome block to another patient with the same number of scans, keeps
features and folds, and re-runs the full nested pipeline 100 times; p = (1 + \#\{perm $\geq$ obs\}) / 101.

\begin{table}[H]
\centering
\small
\begin{tabularx}{\linewidth}{@{}R{4.6cm} L L@{}}
\toprule
 & Grade 4 & All gliomas \\
\midrule
Post-hoc best configuration (optimistic) & boosting, clinical + timing: 0.714 [0.643, 0.785] & ridge, clinical + timing + volumes + morphology + radiomics: 0.744 [0.679, 0.804] \\
Nested over 126 (honest) & 0.664 [0.601, 0.722], slope 0.49 & 0.707 [0.647, 0.766] \\
Permutation null mean (95th percentile) & 0.548 (0.636) & 0.528 \\
Permutation p & 0.050 (4/100 $\geq$ observed) & 0.010 (0/100) \\
\bottomrule
\end{tabularx}
\end{table}

Single groups, grade 4 (ridge / boosting): radiomics 0.675 / 0.628; morphology 0.650 / 0.538; volumes 0.611 /
0.573; clinical 0.585 / 0.612; timing 0.559 / 0.649; change 0.537 / 0.565. The permutation null is centred
above 0.5 because selecting the best of 126 configurations is optimistic even without signal. Boosting was
chosen in 15/25 outer folds; selected boosting models had large train--test gaps (for example training AUC
0.961 for the best grade 4 configuration).

\FloatBarrier
\section{Phase-of-care ladder: full results}\label{sec:S11}

Source: \code{docs/results/PHASE\_LADDER.md};
\code{outputs/results/analyses/phase\_ladder/\{results,cohorts,permutation,
subset\_of\_full,nested\_choices\}.csv};
\code{docs/manuscript/figures/fig4\_phase.csv}.

\textbf{Phase definitions} (grade 4 / all-glioma scans): pre-RT, \code{rt\_started == 0}, including
patients
with no RT documented (76 / 115); during RT, started but not completed (6 / 8); early post-RT, completed and
$\leq$ 84 d after RT end (98 / 104); later post-RT, > 84 d (69 / 81). RT without a documented start date
counts as started; RT without an end date is never ``completed''.

\begin{table}[H]
\centering
\footnotesize
\renewcommand{\thetable}{S15}\renewcommand{\theHtable}{S15}
\caption{\textbf{Cohorts.}}
\label{tab:S15}
\begin{tabularx}{\linewidth}{@{}L R{0.78cm} R{1.24cm} R{0.93cm} R{1.55cm} R{1.24cm} R{2.20cm} R{2.20cm}@{}}
\toprule
Population $\cdot$ cohort & Scans & Patients & Events & Non-events & Censored & EPP R2 (10 params) & EPP R2 + I2 (54) \\
\midrule
Grade 4 $\cdot$ full & 249 & 126 & 93 & 117 & 39 & 9.3 & 1.7 \\
Grade 4 $\cdot$ post-RT & 167 & 84 & 83 & 83 & 1 & 8.3 & 1.5 \\
Grade 4 $\cdot$ later post-RT & 69 & 40 & 36 & 32 & 1 & 3.6 & 0.7 \\
All $\cdot$ full & 308 & 155 & 100 & 156 & 52 & 10.0 & 1.9 \\
All $\cdot$ post-RT & 185 & 92 & 84 & 96 & 5 & 8.4 & 1.6 \\
All $\cdot$ later post-RT & 81 & 46 & 37 & 40 & 4 & 3.7 & 0.7 \\
\bottomrule
\end{tabularx}
\end{table}

By the Riley criteria (anticipated R$^{2}$cs 15\% of maximum), the primary post-RT cohort would need about 750
scans (375 events) for R2 and about 4,040 scans (2,022 events) for R2 + I2. Nearly all censoring is in pre-RT
and early scans. In post-RT cohorts two phase dummies are constant, so R2 has about 8 effective parameters
(exploratory note).

\begin{table}[H]
\centering
\small
\renewcommand{\thetable}{S16}\renewcommand{\theHtable}{S16}
\caption{\textbf{Primary and key secondary results} (grade 4, post-RT, ridge).}
\label{tab:S16}
\begin{tabularx}{\linewidth}{@{}L R{3.42cm} R{4.12cm} R{0.88cm} R{0.80cm}@{}}
\toprule
Model & AUC (95\% CI) & $\Delta$ vs R2 (95\% CI) & Slope & O/E \\
\midrule
R2 (clinical + phase) & 0.549 (0.458--0.642) & --- & 0.76 & 0.99 \\
R2 + I2 (primary) & 0.610 (0.529--0.688) & +0.061 ($-$0.038, +0.156) & 0.68 & 0.99 \\
R2 + I3 (key secondary) & 0.609 (0.543--0.676) & +0.061 ($-$0.041, +0.157) & 0.44 & 0.96 \\
R2 + I1 & 0.593 & +0.044 ($-$0.054, +0.137) & 0.60 & 0.98 \\
R2 + I4 & 0.625 & +0.076 ($-$0.041, +0.192) & 0.30 & 0.97 \\
R3 (baseline), reference & 0.604 (0.520--0.690) & --- & 0.85 & 0.99 \\
\bottomrule
\end{tabularx}
\end{table}

Permutation test of the primary contrast: p = 0.208 (20 of 100 permutations $\geq$ +0.061; p = 21/101); null
mean +0.028, median +0.025, range $-$0.104 to +0.162. The nested radiomics choice was unstable (most common
configuration in 12 of 25 outer folds).

\begin{table}[H]
\centering
\small
\renewcommand{\thetable}{S17}\renewcommand{\theHtable}{S17}
\caption{\textbf{Bases by cohort (grade 4, ridge; AUC [95\% CI])} (\code{fig4\_phase.csv}).}
\label{tab:S17}
\begin{tabularx}{\linewidth}{@{}R{2.02cm} L R{3.60cm} R{3.60cm}@{}}
\toprule
Base & Full & Post-RT & Later post-RT \\
\midrule
R0 clinical & 0.585 [0.496, 0.666] & 0.552 [0.456, 0.646] & 0.343 [0.281, 0.418] \\
R1 + timing & 0.644 [0.560, 0.721] & 0.572 [0.485, 0.667] & 0.388 [0.291, 0.503] \\
R2 + phase & 0.664 [0.569, 0.753] & 0.549 [0.458, 0.642] & 0.376 [0.294, 0.470] \\
R3 baseline & 0.698 [0.615, 0.774] & 0.604 [0.520, 0.690] & 0.472 [0.339, 0.617] \\
\bottomrule
\end{tabularx}
\end{table}

Full grade 4 cohort: R1 $-$ R0 +0.059 ($-$0.016, +0.133); R2 $-$ R0 +0.080 (+0.005, +0.152); R3 $-$ R0 +0.113
(+0.026, +0.208); R3 $-$ R2 +0.034 ($-$0.027, +0.090). R2 + I1 0.696 vs R3 0.698 ($\Delta$ $-$0.002, $-$0.081
to +0.074). Imaging on top of each base (full grade 4; I1 / I2 / I3 / I4): R0 +0.074 / +0.071 / +0.065 /
+0.063; R1 +0.030 / +0.026 / +0.008 / +0.018; R2 +0.031 / +0.023 / +0.002 / +0.020; R3 +0.004 / $-$0.007 /
$-$0.032 / $-$0.004; no CI excludes 0. Post-RT grade 4 (I1 / I2 / I3 / I4): R0 +0.049 / +0.058 / +0.046 /
+0.079; R1 +0.030 / +0.038 / +0.042 / +0.060; R2 +0.044 / +0.061 / +0.061 / +0.076; R3 +0.008 / +0.017 /
+0.005 / +0.039; no CI excludes 0. All gliomas, full: R0 0.619, R1 0.693, R2 0.700, R3 0.720. Boosting
(secondary): grade 4 full R2 0.698, R1 0.693; imaging never helped boosting ($-$0.027 to +0.007); post-RT R2 +
I2 vs R2 +0.012 ($-$0.069, +0.102).

\textbf{Sensitivity analyses.} Later post-RT only (grade 4: 69 scans, 40 patients, 36 events): bases below
chance (R0 0.343, calibration slope $-$2.0), the known negative bias of CV for near-null models with few
events; large $\Delta$AUCs there (e.g. R2 + I3 +0.250) are measured against a below-chance reference and are
not evidence. Full-cohort models evaluated on post-RT scans only: grade 4 R2 0.570, R2 + I2 0.620, $\Delta$
+0.050 ($-$0.061, +0.157); R3 0.638, R3 + I2 $\Delta$ $-$0.014. Exploratory: age alone 0.586 in grade 4
post-RT; age + I2 vs age +0.022 ($-$0.103, +0.145). LUMIERE could not be used: no RT dates; the rater flag
marks 27 of 132 scans and leaves 83 scans with unknown phase.

\FloatBarrier
\section[Secondary analyses: baselines, sensitivity, imbalance, schedule, feature conflicts, SHAP, calibration and
utility]{Secondary analyses}\label{sec:S12}

\subsection{Literature baselines and tuning}\label{sec:S12.1}

Source: \code{docs/results/BASELINES.md}.

\begin{table}[H]
\centering
\footnotesize
\renewcommand{\thetable}{S18}\renewcommand{\theHtable}{S18}
\caption{\textbf{MU ladder (reference = baseline ridge).}}
\label{tab:S18}
\begin{tabularx}{\linewidth}{@{}L R{2.75cm} R{3.5cm} R{2.75cm} R{3.5cm}@{}}
\toprule
Model & Grade 4 AUC [95\% CI] & $\Delta$ [95\% CI] & All AUC [95\% CI] & $\Delta$ [95\% CI] \\
\midrule
Baseline ridge (reference) & 0.698 [0.615, 0.774] & --- & 0.720 [0.638, 0.792] & --- \\
Clinical only & 0.585 [0.496, 0.666] & $-$0.113 [$-$0.208, $-$0.026] & 0.619 [0.540, 0.696] & $-$0.101 [$-$0.177, $-$0.031] \\
Timing only & 0.559 [0.475, 0.637] & $-$0.139 [$-$0.214, $-$0.070] & 0.607 [0.528, 0.684] & $-$0.112 [$-$0.175, $-$0.059] \\
Clinical + timing & 0.644 [0.560, 0.721] & $-$0.054 [$-$0.103, $-$0.012] & 0.693 [0.617, 0.766] & $-$0.027 [$-$0.061, +0.007] \\
B0 log ET volume & 0.636 [0.541, 0.723] & $-$0.062 [$-$0.183, +0.043] & 0.673 [0.591, 0.752] & $-$0.046 [$-$0.142, +0.054] \\
B0 $\Delta$ET vs previous & 0.546 [0.479, 0.614] & $-$0.153 [$-$0.241, $-$0.056] & 0.531 [0.465, 0.590] & $-$0.189 [$-$0.285, $-$0.101] \\
B1 volumetric RANO-style rule & 0.551 [0.480, 0.620] & $-$0.147 [$-$0.239, $-$0.062] & 0.554 [0.488, 0.615] & $-$0.166 [$-$0.258, $-$0.087] \\
B1 rule, RANO 2.0 window & 0.508 [0.440, 0.577] & $-$0.191 [$-$0.289, $-$0.098] & 0.509 [0.448, 0.568] & $-$0.210 [$-$0.301, $-$0.119] \\
B4 Hashim-style Cox & 0.628 [0.538, 0.711] & $-$0.070 [$-$0.179, +0.031] & 0.678 [0.596, 0.755] & $-$0.041 [$-$0.120, +0.039] \\
\bottomrule
\end{tabularx}
\end{table}

The B1 rule: ET $\geq$ 40\% and $\geq$ 1 cc above the past-only nadir (measurable nadir $\geq$ 0.1 cc), or
$\geq$ 1 cc from a $\sim$0 cc nadir; positive in 25\% of grade 4 scans. Nested hyperparameter grids: every
tuned arm within $\pm$0.025 of the baseline, all CIs include 0 (e.g. nested over ridge + boosting $-$0.024
[$-$0.061, +0.011], grade 4); choices were unstable.

\subsection{Sensitivity of the primary comparison}\label{sec:S12.2}

Source: \code{docs/results/SENSITIVITY.md}.

\begin{table}[H]
\centering
\footnotesize
\begin{tabularx}{\linewidth}{@{}L R{1.5cm} R{1.5cm} R{2.9cm} R{1.4cm} R{3.5cm}@{}}
\toprule
Sensitivity & Scans / patients & Events / censored & Baseline AUC [95\% CI] & + radiomics (nested) & $\Delta$AUC [95\% CI] \\
\midrule
Primary: grade 4, 120 d & 249 / 126 & 93 / 39 & 0.698 [0.615, 0.774] & 0.666 & $-$0.032 [$-$0.084, +0.017] \\
Horizon 90 d & 249 / 126 & 62 / 35 & 0.663 [0.577, 0.739] & 0.664 & +0.001 [$-$0.057, +0.063] \\
Horizon 180 d & 249 / 126 & 131 / 43 & 0.667 [0.573, 0.752] & 0.636 & $-$0.032 [$-$0.092, +0.022] \\
Composite (progression or death) & 249 / 126 & 107 / 25 & 0.642 [0.566, 0.713] & 0.629 & $-$0.013 [$-$0.073, +0.052] \\
Alternative progression date & 244 / 124 & 92 / 39 & 0.670 [0.578, 0.757] & 0.648 & $-$0.023 [$-$0.066, +0.019] \\
Imaging-only follow-up & 249 / 126 & 93 / 43 & 0.685 [0.600, 0.761] & 0.667 & $-$0.018 [$-$0.070, +0.033] \\
All gliomas & 308 / 155 & 100 / 52 & 0.720 [0.638, 0.792] & 0.747 & +0.028 [$-$0.017, +0.071] \\
Excluding 12-week post-RT scans & 151 / 90 & 46 / 39 & 0.725 [0.602, 0.834] & 0.727 & +0.002 [$-$0.064, +0.067] \\
Excluding patient 0145 (conflicting dates) & 247 / 125 & 93 / 39 & 0.691 [0.602, 0.776] & 0.687 & $-$0.004 [$-$0.058, +0.048] \\
\bottomrule
\end{tabularx}
\end{table}

\subsection{Class imbalance}\label{sec:S12.3}

Source: \code{docs/results/IMBALANCE.md}. Event shares are 44\% (grade 4) and 39\% (all
gliomas), a mild
imbalance. Over learners (ridge, boosting, XGBoost, LightGBM, nested variants) $\times$ re-balancing (none,
class weights, patient-aware oversampling, SMOTE with neighbours from other patients), $\Delta$AUC against the
ridge baseline ranged from $-$0.031 to +0.023; 4 of 190 CIs excluded 0 (ridge, all gliomas, +0.009 to +0.014).
Re-balancing shifted predicted risk upward as expected from a prior shift (median O/E 1.01 without
re-balancing vs 0.88--0.89 with it) and raised Brier. Grade 4 ridge with class weights: $\Delta$AUC +0.008
[$-$0.001, +0.019], O/E 1.04 $\rightarrow$ 0.94. Fixed small-data XGBoost: +0.006 [$-$0.060, +0.072] (grade
4); nested-tuned XGBoost had calibration slope about 0.6 and transported worse to LUMIERE (0.562 vs 0.608). No
re-balancing is used in any reported model.

\subsection{Schedule coupling}\label{sec:S12.4}

Source: \code{docs/results/SCHEDULE\_COUPLING.md}. The interval to the next
MRI is unknown at prediction time
and partly caused by the outcome (scans are brought forward when progression is suspected), so conditioning on
it is a conservative test. Grade 4: positive rate 0.54 when the next scan is $\leq$ 60 d away and 0.10 when >
180 d; interval alone AUC 0.639 [0.555, 0.721]. Baseline AUC within interval tertiles 0.718 / 0.762 / 0.617;
conditional AUC (pairs within a tertile) 0.687 [0.595, 0.768]; logit coefficient per SD 0.91 $\rightarrow$
0.86 [0.43, 1.41] after adjusting for log interval; baseline features cannot predict the interval (grouped-CV
R$^{2}$ $-$0.05). Excluding early post-RT scans and scans whose next scan is > 120 d away: 0.777 [0.653,
0.891] (85 scans). Nested radiomics conditional AUC 0.654, below the baseline. In LUMIERE the interval alone
scores 0.788 [0.703, 0.868], and the + volumes boosting model's coefficient is no longer distinguishable from
0 after adjustment (0.37 [$-$0.14, 0.89]).

\subsection{Feature conflicts}\label{sec:S12.5}

Source: \code{docs/results/FEATURE\_CONFLICTS.md}. No feature group
measurably hurts: every
leave-one-group-out contrast that would mean ``group hurts'' has a CI including 0; adding any imaging group to
the baseline changes AUC by $-$0.004 to +0.019 (CIs span 0). The conflict is redundancy: timing and treatment
encode the same RT-status information (interaction $-$0.066 to $-$0.073, CI excluding 0 in 3 of 4 population
$\times$ learner cells). Within the baseline, 9 design columns are exactly collinear (VIF = $\infty$) and 16
have VIF > 10; ``RT completed'' appears in at least five columns. Pre-stated pruning evaluated nested:
$\Delta$AUC $-$0.018 to +0.011, all CIs include 0; the same rule applied once globally gives +0.009 to +0.018,
which is selection optimism.

\subsection{SHAP attributions}\label{sec:S12.6}

Source: \code{docs/results/SHAP.md}. Ridge: exact linear SHAP on in-fold standardised
features; boosting:
TreeSHAP (path-dependent), additivity checked to $10^{-6}$; explanations out of fold, 25 fits. Group share =
mean |per-scan group sum| / sum over groups.

\begin{table}[H]
\centering
\small
\begin{tabularx}{\linewidth}{@{}L R{1.2cm} R{6.5cm}@{}}
\toprule
Model (grade 4) & OOF AUC & Group shares \\
\midrule
Baseline ridge & 0.698 & treatment 58\%, clinical 25\%, timing 17\% \\
Baseline boosting & 0.687 & treatment 51\%, clinical 35\%, timing 14\% \\
Full non-radiomics ridge & 0.692 & treatment 34\%, morphology 21\%, change 15\%, clinical 13\%, volumes 10\%, timing 8\% \\
Full non-radiomics boosting & 0.694 & treatment 26\%, morphology 23\%, clinical 21\%, change 13\%, volumes 13\%, timing 3\% \\
Baseline + radiomics score (post hoc configuration) & 0.721 & treatment 48\%, clinical 20\%, timing 16\%, radiomics 16\% \\
\bottomrule
\end{tabularx}
\end{table}

Time-relative features hold 35\% (ridge) / 47\% (boosting) of feature-level credit in the baseline. Imaging
takes 45\% / 51\% of credit in the full models without an AUC gain. On LUMIERE, the locked ridge's
under-prediction is driven by \code{idh\_unknown} (mean $-$0.71 logits; mean |SHAP| 0.77 on LUMIERE vs
0.13 on
MU), by lower DeepBraTumIA ET volumes ($-$0.32 logits) and by a missingness indicator that is never set in
LUMIERE.

\subsection{Calibration, decision curves, added value and sample size}\label{sec:S12.7}

Source: \code{docs/results/CLINICAL\_UTILITY.md}.

\begin{table}[H]
\centering
\small
\begin{tabularx}{\linewidth}{@{}L R{0.97cm} R{2.20cm} R{3.42cm} R{3.07cm} R{0.97cm}@{}}
\toprule
Model (grade 4) & AUC & O/E (weighted) & CITL [95\% CI] & Slope [95\% CI] & ICI \\
\midrule
Baseline ridge & 0.698 & 0.96 & $-$0.09 [$-$0.42, 0.28] & 0.89 [0.50, 1.43] & 0.032 \\
Baseline boosting & 0.687 & 0.94 & $-$0.12 [$-$0.47, 0.23] & 0.82 [0.44, 1.31] & 0.039 \\
Baseline + radiomics (nested) & 0.666 & 0.94 & $-$0.13 [$-$0.46, 0.25] & 0.60 [0.33, 0.93] & 0.063 \\
Baseline + deep (nested D1) & 0.694 & 0.94 & $-$0.15 [$-$0.52, 0.28] & 0.55 [0.28, 0.91] & 0.085 \\
\bottomrule
\end{tabularx}
\end{table}

\begin{table}[H]
\centering
\small
\begin{tabularx}{\linewidth}{@{}L R{3.5cm} R{3.5cm} R{3.5cm}@{}}
\toprule
Net benefit (grade 4, IPCW) & 10\% & 20\% & 30\% \\
\midrule
Treat all & 0.369 & 0.291 & 0.189 \\
Baseline ridge & 0.372 & 0.300 & 0.232 \\
$\Delta$ radiomics (nested) & $-$0.000 [$-$0.004, 0.004] & $-$0.001 [$-$0.013, 0.011] & $-$0.020 [$-$0.038, $-$0.004] \\
$\Delta$ boosting & $-$0.002 [$-$0.007, 0.001] & $-$0.007 [$-$0.021, 0.006] & $-$0.011 [$-$0.039, 0.013] \\
$\Delta$ deep (nested D1) & $-$0.002 [$-$0.013, 0.006] & +0.002 [$-$0.014, 0.018] & +0.001 [$-$0.023, 0.026] \\
\bottomrule
\end{tabularx}
\end{table}

Added value vs baseline (continuous NRI; IDI): radiomics +0.17 [$-$0.05, 0.41]; +0.000 [$-$0.030, 0.029]. Deep
+0.48 [0.24, 0.73]; +0.036 [$-$0.009, 0.082]. The positive NRI reflects upward risk shifts in events by
over-fitted models; $\Delta$AUC and $\Delta$ net benefit are null \citep{pencina2008}. All gliomas: radiomics
$\Delta$ net benefit +0.011 [0.004, 0.020] at 10\% and +0.014 [0.001, 0.027] at 20\%, with $\Delta$AUC +0.028
[$-$0.017, +0.071]; nothing similar in grade 4. LUMIERE locked boosting: net benefit $\approx$ treat-all
(0.218 vs 0.222 at 20\%).

\textbf{Sample size (\citealp{riley2020}; grade 4).} 210 verified scans, 93 events, outcome proportion 0.443,
R$^{2}$cs 0.119 (from AUC 0.698). Baseline with 31 design columns: 3.0 events per parameter; required 2,190
scans (970 events) by the shrinkage criterion. Adding 10 radiomics features: 2,896; the full radiomics pool
(2,477 candidate parameters): 174,942. Requirements assume independent scans, so the shortfall is understated.

\FloatBarrier
\section{External evaluation on LUMIERE: full results}\label{sec:S13}

Source: \code{docs/results/EXTERNAL\_LUMIERE.md};
\code{outputs/results/external\_lumiere/results.csv}.

Shared features: age, sex, IDH and MGMT (with unknown), log days since surgery, < 3 months post-RT flag, log
ET / necrosis + cavity / SNFH volumes (DeepBraTumIA; 14 of 132 index scans without usable volumes, imputed
with MU training medians inside the locked pipeline).

\begin{landscape}
\begin{table}[H]
\centering
\small
\renewcommand{\thetable}{S19}\renewcommand{\theHtable}{S19}
\caption{\textbf{Models trained on MU grade 4.} Pooled locked AUC [95\% CI] and O/E; other columns are per-fold means.}
\label{tab:S19}
\begin{tabularx}{\linewidth}{@{}L R{1.6cm} R{1.2cm} R{3.4cm} R{1.0cm} R{1.5cm} R{1.6cm} R{1.6cm} R{1.6cm}@{}}
\toprule
Features & Learner & MU CV & Locked, pooled & O/E & Locked, per fold & Fine-tuned (shrink) & Fine-tuned (pooled) & LUMIERE only \\
\midrule
Clinical & ridge & 0.585 & 0.610 [0.487, 0.737] & 1.15 & 0.558 & 0.544 & 0.505 & 0.396 \\
Clinical & boosting & 0.625 & 0.627 [0.514, 0.731] & 0.86 & 0.624 & --- & 0.560 & 0.520 \\
Clinical + timing & ridge & 0.612 & 0.610 [0.497, 0.737] & 0.98 & 0.603 & 0.591 & 0.603 & 0.586 \\
Clinical + timing & boosting & 0.664 & 0.632 [0.501, 0.748] & 0.84 & 0.643 & --- & 0.632 & 0.627 \\
Volumes & ridge & 0.611 & 0.545 [0.410, 0.653] & 1.00 & 0.521 & 0.531 & 0.567 & 0.569 \\
Volumes & boosting & 0.591 & 0.529 [0.389, 0.644] & 1.00 & 0.515 & --- & 0.513 & 0.518 \\
\textbf{Clinical + timing + volumes} & \textbf{ridge} & 0.667 & \textbf{0.608 [0.469, 0.739]} & \textbf{1.33} & 0.598 & 0.606 & 0.617 & 0.580 \\
\textbf{Clinical + timing + volumes} & \textbf{boosting} & 0.675 & \textbf{0.620 [0.495, 0.733]} & \textbf{0.99} & 0.608 & --- & 0.658 & 0.610 \\
\bottomrule
\end{tabularx}
\end{table}
\end{landscape}

Recalibration equals locked per fold (pooled O/E 0.93--0.97). No adaptation arm has a $\Delta$AUC CI above 0;
LUMIERE-only is significantly worse for clinical ridge ($\Delta$ $-$0.162 [$-$0.238, $-$0.066]). Cohort
normalisation of volumes (rank) raised the pooled locked AUC of every volume model (+0.007 to +0.069);
clinical + timing + volumes boosting 0.620 $\rightarrow$ 0.659 [0.529, 0.767], per-fold $\Delta$ +0.047
[$-$0.002, +0.110]. This is transductive and a sensitivity analysis only. All-glioma-trained locked clinical +
timing + volumes: ridge 0.599 [0.462, 0.728] (O/E 1.30), boosting 0.619 [0.477, 0.740] (O/E 1.04).

\textbf{Deep features on LUMIERE} (118 scans with deep features;
\code{deep/lumiere/locked.csv}): grade 4
clinical + timing 0.622 $\rightarrow$ + deep 0.679 ($\Delta$ +0.057 [$-$0.132, +0.235], O/E 2.08); random
encoder same configuration 0.667 (+0.045). With volumes: 0.635 $\rightarrow$ 0.725 (+0.089 [$-$0.026,
+0.215]). All gliomas: +0.064 [$-$0.011, +0.144]; with volumes +0.050 [+0.002, +0.109]. Pretrained vs random,
same configuration: +0.011 to +0.069, no CI above 0. A classifier separates LUMIERE from MU on deep features
with AUC 0.82 / 0.79 / 0.81 (BraTS 2021 / BrainSegFounder / random).

\textbf{PRISM's initial radiomics model, as shipped} (legacy features re-extracted, 118 scans; files dated 2026-10-09, before
the final LUMIERE label definition): raw 0.613 [0.500, 0.729] pooled (O/E 1.24); recalibrated O/E 0.98;
fine-tuned 0.570 pooled.

\textbf{Literature baselines on LUMIERE} (locked; \code{BASELINES.md}): clinical + timing 0.610 [0.497,
0.737]; Hashim-style Cox 0.640 [0.508, 0.754] ($\Delta$ +0.030 [$-$0.080, +0.127]); volumetric RANO-style rule
0.559.

\textbf{Harmonisation note.} DeepBraTumIA ET and SNFH medians (grade 4: 0.61 and 19.11 cc) are far below MU's
(5.45 and 42.85 cc). The harmonised necrosis + cavity definition is correct but did not improve transfer in a
development run (\code{FIX\_FEATURES.md} \S{}1).

\FloatBarrier
\section{Imaging-defined label, imaging diagnostics and re-segmentation}\label{sec:S14}

Source: \code{docs/results/CLEAN\_LABEL.md};
\code{docs/results/DIAGNOSTICS\_IMAGING.md};
\code{docs/results/RESEGMENTATION.md};
\code{outputs/results/analyses/\{clean\_label,diagnostics\_imaging,reseg\}/}.

\textbf{Imaging-defined label.} MU codes the type of first progression: imaging, clinical only, or both.
Censoring the clinical-only progressions at their progression date changed grade 4 events from 93 (44 / 12 /
37 by type) to 81 and censored scans from 39 to 51; no non-event changed. Primary (grade 4, full, ridge):
clinical 0.563 (0.474--0.642); baseline 0.696 (0.610--0.776); imaging (volumes + change) 0.628 (0.549--0.701),
$\Delta$ $-$0.068 ($-$0.167, +0.033); baseline + imaging 0.722 (0.639--0.799), $\Delta$ +0.027 ($-$0.033,
+0.082). On the original label, baseline + imaging was 0.694 ($\Delta$ $-$0.004, $-$0.060 to +0.051).
Secondary contrasts (post-RT, all gliomas, boosting) gave baseline + imaging $\Delta$ +0.008 to +0.028 with
all CIs including 0; imaging alone in all gliomas was worse than the baseline ($-$0.090, $-$0.176 to
$-$0.007). Within-LUMIERE CV with the expert RANO label: every model at or below chance (ridge clinical +
timing 0.466), so LUMIERE cannot answer the label question. Rule not met.

\textbf{Volume cross-check against the dataset's own volume table (pass).} Pearson r: ET 0.9998 (exact voxel
match 98.3\%), SNFH 0.9994 (99.7\%), RC 0.9993 (94.1\%), NETC 0.9957 (83.7\%); disagreements are concentrated
in seven patients with extra timepoints or a second mask version in the table. Measurement volumes equal mask
voxel counts exactly. Label mapping: diagonal agreement 0.94--0.97, off-diagonal $\leq$ 0.003. Intensities
agree with labels (ET/brain on T1c median 1.45, > 1 in 98\% of scans; SNFH/brain on FLAIR median 1.44, > 1 in
100\%).

\textbf{Concurrent detection (positive control; fail).} Progression scan = first MRI with a mask on or after
the documented progression date (120 scans, all gliomas).

\begin{table}[H]
\centering
\small
\begin{tabularx}{\linewidth}{@{}L R{4.30cm} R{2.02cm}@{}}
\toprule
Comparison (imaging only, 5 $\times$ 5 grouped CV; ridge / boosting) & All gliomas & Grade 4 \\
\midrule
Progression scans vs all index scans & 0.62 [0.57, 0.67] / 0.63 & 0.62 / 0.63 \\
Progression scans vs verified 120-d negatives & 0.68 [0.62, 0.74] / 0.70 & 0.65 / 0.66 \\
Both groups with a prior scan & 0.62 / 0.61 & 0.60 / 0.62 \\
\bottomrule
\end{tabularx}
\end{table}

Progression scan vs the same patient's previous scan (74 pairs): ET larger in 49\% (Wilcoxon p = 0.54); lesion
larger in 65\% (p = 0.033); lesion growth from nadir larger in 79\% (n = 42, p = 0.0003); ET growth from nadir
50\% (p = 0.95). By progression type: imaging (40 pairs) ET larger in 50\%; imaging + clinical (27) 52\%;
clinical only (7) 29\%.

\textbf{Why the control fails.} Between consecutive scans $\leq$ 14 days apart (n = 25), the median change was
3.2-fold in ET and 1.8-fold in lesion volume, larger than over 45--90 day gaps; across 389 pairs the size of
the change was unrelated to the gap (Spearman $-$0.01, p = 0.83). Some short-gap jumps straddle a
re-resection; others (four patients) are implausible.

\textbf{Label--date alignment (pass).} The progression date coincides with an MRI for 62\% of progressors
($\pm$7 d: 69\%; 30 d: 73\%); 22\% of progressors have no MRI on or after the progression date. Positives:
median time to progression 66 d (IQR 44--92); 76\% are the last scan before progression. Agreement between the
120-d label and ``the next scan is the progression scan'': 79\% ($\kappa$ = 0.53).

\textbf{Horizon curve (grade 4, IPCW AUC(H); imaging / clinical + phase / both).} 60 d: 0.63 / 0.55 / 0.65; 90
d: 0.62 / 0.63 / 0.68; 120 d: 0.60 / 0.66 / 0.69; 180 d: 0.58 / 0.59 / 0.63; 270 d: 0.59 / 0.55 / 0.63.
Imaging stays at about 0.6 at every horizon; 30-d estimates (16 events) are not interpretable.

\textbf{Data issues to report to the data owners.} Three scan pairs from different patients have identical
masks (0018\_T5 / 0032\_T5, 0060\_T2 / 0062\_T2, 0113\_T1 / 0132\_T1) and four further pairs are
near-identical; a few implausible within-days volume jumps; 15 of 96 recorded deaths precede a dated MRI (S3).

\textbf{Re-segmentation.} BraTS 2021 SwinUNETR and BrainSegFounder were run on all 594 complete scans (sliding
window, ROI 128$^{3}$, overlap 0.5; labels written in the MU convention; RC taken from the MU mask where the
new map is background, because BraTS models have no cavity class and label most of the cavity as NETC, so
tumour core is not comparable).

\begin{table}[H]
\centering
\small
\begin{tabularx}{\linewidth}{@{}L R{3.07cm} R{3.07cm} R{2.20cm}@{}}
\toprule
Agreement with MU masks & Dice ET median [IQR] & Dice WT & Pearson log-volume ET / WT \\
\midrule
BraTS 2021 SwinUNETR & 0.80 [0.49, 0.91] & 0.85 [0.74, 0.93] & 0.92 / 0.91 \\
BrainSegFounder & 0.81 [0.56, 0.91] & 0.86 & 0.94 / 0.92 \\
\bottomrule
\end{tabularx}
\end{table}

Short-gap consistency ($\leq$ 30 d, 54 pairs): median ET fold change 2.57 $\rightarrow$ 2.17 ($\Delta$ median
|log ratio| $-$0.17 [$-$0.46, $-$0.02]); WT unchanged. Positive controls with the new masks: primary detection
(progression scan vs verified negatives, ridge) 0.681 $\rightarrow$ 0.665, $\Delta$ $-$0.016 [$-$0.057,
+0.025] (BrainSegFounder 0.677); ET larger at progression 49\% $\rightarrow$ 45\% (BrainSegFounder 41\%).
Forward 120-d task, imaging only: all gliomas 0.644 $\rightarrow$ 0.675, $\Delta$ +0.031 [$-$0.017, +0.084];
clinical + phase + imaging 0.712 $\rightarrow$ 0.740, $\Delta$ +0.029 [$-$0.010, +0.069]. An exploratory
plain-AUC boosting gain (+0.059 [+0.014, +0.107]) was one of more than 20 comparisons, not pre-specified and
not confirmed by the IPCW ridge estimate. Verdict by the pre-specified rule: the new segmentation does not
help.

\FloatBarrier
\section{The auditable layer: full results and an example audit record}\label{sec:S15}

Source: \code{docs/results/AUDIT\_SYSTEM.md};
\code{outputs/results/analyses/audit\_system/}.
Run: 2026-10-15,
\code{python -m prism audit --workers 24}, 114 s. Checks (\code{checks.csv}): CV predictions reproduce the
baseline OOF predictions (max |$\Delta$| \num{1.7e-16}); intercept + $\Sigma$ contributions = logit (max
error \num{1.3e-15}). Code changes after the plan were confined to an as-of rule for demographics, a NaN
rule for one-class strata, and a shared estimator for referral subsets; none changed a model or a rule.

\textbf{Locked artifact and model card.} Grade 4: 249 training scans, 126 patients, 210 verified, 93 events; C
= 0.1; conformal threshold q at $\alpha$ = 0.2 / 0.1: 0.613 / 0.683 (all gliomas 0.529 / 0.620). Out of scope:
pseudoprogression diagnosis; cohorts without treatment fields (including LUMIERE); other institutions without
validation.

\begin{table}[H]
\centering
\small
\renewcommand{\thetable}{S20}\renewcommand{\theHtable}{S20}
\caption{\textbf{Example audit records} (out-of-fold, repeat 0, grade 4; label shown for evaluation only).}
\label{tab:S20}
\begin{tabularx}{\linewidth}{@{}R{1.3cm} R{2.2cm} R{2.6cm} R{1.6cm} L R{1.9cm}@{}}
\toprule
Scan & Phase & Risk [95\% bootstrap] & Conformal set $\alpha$ = 0.2 & Flags & 120-d label \\
\midrule
0201\_T3 & pre-RT & 0.04 [0.01, 0.27] & \{0\} & OOD (MGMT unknown) & 0 \\
0136\_T2 & during RT & 0.47 [0.36, 0.77] & \{0\} & none & 1 \\
0117\_T1 & early post-RT & 0.50 [0.39, 0.61] & \{0, 1\} & (i) 26 d after RT; (iii) & 1 \\
0059\_T1 & early post-RT & 0.48 [0.22, 0.77] & \{0, 1\} & (i); (ii) rt\_start\_day below training 1st percentile; (iii) & 0 \\
0089\_T5 & later post-RT & 0.45 [0.34, 0.72] & \{0, 1\} & (ii) rt\_start\_day above 99th percentile; (iii) & 1 \\
0014\_T4 & later post-RT & 0.63 [0.44, 0.86] & \{1\} & none & 0 (set misses) \\
\bottomrule
\end{tabularx}
\end{table}

\textbf{Example record (PatientID\_0117\_Timepoint\_1, day 85 from diagnosis), abridged from
\code{records.md}.}

\begin{itemize}
\item Phase of care: early post-RT ($\leq$ 84 d after RT end).
\item Risk of documented progression within 120 d: 0.50 (95\% bootstrap interval 0.39--0.61); conformal set at
$\alpha$ = 0.2 \{0, 1\}.
\item Referral: REFER --- early post-RT window (scan 26 d after RT end); uncertain (the $\alpha$ = 0.2 set
contains both classes).
\item Log-odds: intercept $-$0.33 + clinical +0.22 + timing $-$0.14 + treatment +0.24 = $-$0.01 (check error
\num{5.6e-17}).
\item Inputs known by scan day: yes, all. Examples: \code{rt\_started} = 1 (source event day 21 $\leq$ scan
day); \code{rt\_completed} = 1 (day 59);
\code{days\_since\_rt\_end} = 26; prescription 60 Gy in 30 fractions
(visible from day 21); \code{surgery\_day} $-$4.
\item Largest terms: \code{early\_post\_rt} $-$0.12; \code{visit\_index} $-$0.12;
\code{previous\_brain\_tumor} +0.07; \code{rt\_started} +0.07; \code{male}
$-$0.07; \code{mgmt\_unknown} +0.07
(log-odds).
\item Evaluation only: 120-d label 1 (progression at 111 d).
\end{itemize}

The displayed risk (model fitted on the whole training fold) and the conformal set (model fitted on 70\% of
training patients) can disagree in direction (e.g. 0136\_T2: risk 0.47, set \{0\}); this occurred in 1.8\% /
2.9\% of singleton sets (exploratory). Median width of the 95\% bootstrap interval: 0.37 (IQR 0.32--0.44) in
grade 4 and 0.31 in all gliomas.

\begin{table}[H]
\centering
\footnotesize
\renewcommand{\thetable}{S21}\renewcommand{\theHtable}{S21}
\caption{\textbf{P1, conformal coverage at $\alpha$ = 0.2} (\code{conformal.csv}).}
\label{tab:S21}
\begin{tabularx}{\linewidth}{@{}R{1.6cm} R{2.9cm} L R{1.5cm} R{1.3cm} R{1.3cm} R{1.6cm} R{1.3cm}@{}}
\toprule
Population & Overall [95\% CI] & Pre-RT & During RT & Early post-RT & Later post-RT & Events / non-events & Singleton rate \\
\midrule
Grade 4 & 0.841 [0.790, 0.885] & 0.868 (0/5 events covered) & 0.567 (6 scans) & 0.871 & 0.806 & 0.824 / 0.855 & 0.595 \\
All gliomas & 0.823 [0.776, 0.865] & 0.850 (0/10) & 0.550 (8 scans) & 0.852 & 0.790 & 0.698 / 0.904 & 0.696 \\
\bottomrule
\end{tabularx}
\end{table}

Per-patient coverage 0.828 / 0.796. At $\alpha$ = 0.1: 0.949 [0.916, 0.975] / 0.928 [0.896, 0.955], singleton
rate 0.296 / 0.422. Conformal model AUC 0.667 / 0.709 (cost of the 70/30 split). Displayed-risk AUC on
singleton scans 0.732 / 0.723 (non-singletons 0.628 / 0.594).

\begin{table}[H]
\centering
\footnotesize
\renewcommand{\thetable}{S22}\renewcommand{\theHtable}{S22}
\caption{\textbf{P2 and referral rules} (\code{p2.csv}, \code{referral.csv},
\code{referral\_within\_phase.csv}).}
\label{tab:S22}
\begin{tabularx}{\linewidth}{@{}R{1.7cm} R{1.5cm} L R{1.9cm} R{1.5cm} R{3.6cm}@{}}
\toprule
 & Realised coverage & AUC retained [95\% CI] & Composition-matched & Random removal & P2 [95\% CI] \\
\midrule
Grade 4 & 0.671 & 0.731 [0.635, 0.812] & 0.725 & 0.698 & +0.006 [$-$0.023, +0.036] \\
All gliomas & 0.685 & 0.735 [0.654, 0.820] & 0.733 & 0.720 & +0.003 [$-$0.031, +0.036] \\
\bottomrule
\end{tabularx}
\end{table}

\begin{table}[H]
\centering
\footnotesize
\begin{tabularx}{\linewidth}{@{}R{2.4cm} L R{1.3cm} R{1.5cm} R{2.6cm} R{2.6cm}@{}}
\toprule
Grade 4 rule & Referred & AUC kept & AUC referred & Random removal & Composition-matched \\
\midrule
(i) early post-RT & 39.4\% (98 scans, 47 events) & 0.723 & 0.663 & 0.699 [0.673, 0.723] & 0.723 (the rule itself) \\
(ii) out of distribution & 29.8\% & 0.671 & 0.665 & 0.698 & 0.695 \\
(iii) uncertain & 40.5\% & 0.732 & 0.628 & 0.699 & 0.734 [0.709, 0.760] \\
any of (i)--(iii) & 72.6\% & 0.735 & 0.687 & 0.698 & 0.753 \\
\bottomrule
\end{tabularx}
\end{table}

Within-phase AUC (grade 4 / all): pre-RT 0.555 / 0.493; during RT 0.448 / 0.671; early post-RT 0.663 / 0.678;
later post-RT 0.597 / 0.644. On OOD-referred scans the model over-predicts (O/E 0.70, slope 0.66).

\textbf{Shortcut audit.} Unknown marker share by 120-d label (progression / no progression / censored), grade
4: IDH unknown 0/93, 6/117, 5/39; MGMT unknown 6/93, 25/117, 8/39. Dropping both unknown indicators: grade 4
0.676 [0.592, 0.752], $\Delta$ $-$0.023 [$-$0.047, $-$0.001]; also dropping
\code{rt\_date\_unknown} $\Delta$
$-$0.022 [$-$0.045, $-$0.001]; all gliomas $\Delta$ $-$0.011 [$-$0.031, +0.008]. Timing only 0.559, $\Delta$
$-$0.139 [$-$0.214, $-$0.070]; treatment only 0.701 [0.622, 0.776], $\Delta$ +0.003 [$-$0.050, +0.048] (all
gliomas 0.710, $\Delta$ $-$0.009).

\begin{table}[H]
\centering
\small
\renewcommand{\thetable}{S23}\renewcommand{\theHtable}{S23}
\caption{\textbf{Subgroups (grade 4; \textdagger{} fewer than 10 events, descriptive).}}
\label{tab:S23}
\begin{tabularx}{\linewidth}{@{}R{3.3cm} R{2.6cm} L R{2.0cm} R{2.7cm}@{}}
\toprule
Stratum & Scans / events & AUC [95\% CI] & Slope & O/E [95\% CI] \\
\midrule
GBM & 214 / 86 & 0.68 [0.58, 0.76] & 0.78 & 1.10 [0.90, 1.31] \\
Non-GBM grade 4 \textdagger{} & 35 / 7 & 0.72 [0.54, 0.95] & 1.34 & 0.66 [0.30, 1.00] \\
MGMT methylated & 88 / 34 & 0.64 [0.52, 0.77] & 0.62 & 0.99 \\
MGMT unmethylated & 122 / 53 & 0.72 [0.61, 0.84] & 1.11 & 1.14 \\
IDH wild type & 208 / 87 & 0.66 [0.56, 0.77] & 0.78 & 1.09 \\
Male / female & 119 / 44, 130 / 49 & 0.67 [0.55, 0.79] / 0.72 [0.61, 0.83] & 0.85 / 0.94 & 1.02 / 1.06 \\
Age $\leq$ 56.7 / (56.7, 67] / > 67 & 109 / 30, 68 / 35, 72 / 28 & 0.71 / 0.64 / 0.68 & 0.89 / 0.73 / 0.89 & 0.81 / 1.39 [1.07, 1.77] / 1.06 \\
Early post-RT & 98 / 47 & 0.66 [0.55, 0.76] & 1.06 & 0.95 \\
Later post-RT & 69 / 36 & 0.60 [0.45, 0.74] & 0.50 & 1.03 \\
\bottomrule
\end{tabularx}
\end{table}

The middle age tertile is under-predicted (O/E CI excludes 1); uncorrected for 19 strata. Field strength,
scanner and race could not be audited (no acquisition metadata in MU).

\textbf{LUMIERE transport audit} (shared-feature ridge, grade-4-trained; reproduces AUC 0.608, O/E 1.33).
Largest standardised mean differences MU $\rightarrow$ LUMIERE: \code{idh\_unknown} +0.97 (4.4\%
$\rightarrow$
40.9\%), \code{log\_et\_cc} $-$0.84, \code{early\_post\_rt}
$-$0.72, \code{log\_snfh\_cc} $-$0.56. Mean
contribution shift: \code{idh\_unknown} $-$0.68 logits; net linear-predictor shift $-$0.64. Refit
without the
molecular unknown indicators: LUMIERE O/E 1.33 $\rightarrow$ 0.89, slope 0.38 $\rightarrow$ 0.73, AUC 0.608
$\rightarrow$ 0.614 (all gliomas O/E 1.30 $\rightarrow$ 0.96). MU-calibrated conformal coverage on LUMIERE
($\alpha$ = 0.2): 0.715 (events 0.340, non-events 0.947).

\begin{table}[H]
\centering
\small
\renewcommand{\thetable}{S24}\renewcommand{\theHtable}{S24}
\caption{\textbf{Explanation stability} (\code{stability\_summary.csv}).}
\label{tab:S24}
\begin{tabularx}{\linewidth}{@{}L R{2.3cm} R{2.3cm} R{2.3cm} R{2.3cm}@{}}
\toprule
 & Grade 4: 25 CV fits & Grade 4: 200 bootstraps & All: CV & All: bootstrap \\
\midrule
Kendall $\tau$-b, |coef| vs locked, median (IQR) & 0.63 (0.56--0.67) & 0.42 (0.36--0.49) & 0.54 (0.48--0.60) & 0.39 (0.32--0.46) \\
Top-5 overlap, mean (all 5 identical) & 0.67 (4\%) & 0.52 (1.5\%) & 0.65 (4\%) & 0.52 (2\%) \\
Minimum sign agreement over the top 5 & 1.00 & 0.92 & 1.00 & 0.97 \\
Features with sign agreement $\geq$ 0.9 & 61\% & 55\% & 68\% & 61\% \\
\bottomrule
\end{tabularx}
\end{table}

Locked top 5 (grade 4): \code{rt\_start\_day},
\code{previous\_brain\_tumor}, \code{idh\_unknown},
\code{rt\_started}, \code{biopsy\_before\_resection}. Group shares of mean
|contribution|: treatment 51\%,
timing 25\%, clinical 24\%. On a LUMIERE refit of the shared-feature model, 9 of 11 coefficients keep the MU
sign; \code{early\_post\_rt} and the necrosis + cavity volume flip.

\FloatBarrier
\section{Reproducibility}\label{sec:S16}

Every stage is a command of one Python package (\code{python -m prism <stage>}), wrapped by numbered scripts
in \code{scripts/} (\code{run\_all.sh} with a \code{--to STEP} option). Shared fold
assignments are in
\code{outputs/data/folds.csv}; every analysis writes its outputs, a
\code{run\_info.json} and a log under
\code{outputs/}. Seeds are fixed (bootstrap, permutations, conformal splits). The audit layer writes a
\code{MANIFEST.json} with SHA-256 hashes of inputs, code and outputs, the git state and package versions (one
recorded drift: scipy 1.17.1 installed against 1.18.1 in the lock file). Pretrained weights are downloaded
locally and not redistributed. Unit tests cover label rules, as-of-scan treatment features, feature fixes and
the repeated-CV estimator. Commands for each analysis are listed at the end of the corresponding result
document in \code{docs/results/}. \flag{[PENDING: public code release and tag.]}

\FloatBarrier
\section[Reporting checklists (TRIPOD+AI, CLEAR) cross-referenced to the manuscript]{Reporting checklists
cross-referenced to the manuscript}\label{sec:S17}

Item wording: TRIPOD+AI \citep{collins2024},
\code{docs/manuscript/CHECKLIST\_TRIPOD\_AI.md}; CLEAR
(\citealp{kocak2023} \flag{[CHECK]}),
\code{docs/manuscript/CHECKLIST\_CLEAR.md}. ``M'' = main text section;
``S'' = Supplement section.

\clearpage
\textbf{TRIPOD+AI.}

\begin{table}[H]
\centering
\footnotesize
\begin{tabularx}{\linewidth}{@{}R{1.1cm} R{5.4cm} L@{}}
\toprule
Item & Topic & Where reported \\
\midrule
1 & Title & Title (development + external evaluation, population, outcome) \\
2 & Abstract & Structured abstract \\
3a--c & Background, intended use, inequalities & M1; intended use in Key points and S15 model card; inequalities: race not auditable (M4 Limitations, S15) \\
4 & Objectives & M1 final paragraph \\
5a--b & Data sources, dates & M2.3; accrual dates: \flag{[PENDING --- relative days only in both datasets; take from dataset papers]} \\
6a--c & Setting, eligibility, treatments & M2.3--2.4, S2; as-of treatment handling M2.5--2.6, S4 \\
7 & Data preparation & M2.5, S4, S5, S14 \\
8a--c & Outcome, assessors, blinding & M2.4, S3; MU assessor not documented, LUMIERE expert RANO; no blinding (M4 Limitations) \\
9a--c & Predictors & M2.5, S4--S6; MU masks from the dataset, LUMIERE automatic \\
10 & Sample size & Fixed by public datasets; Riley calculation S12.7, M3.8 \\
11 & Missing data & In-fold imputation with indicators (M2.5, S4) \\
12a--g & Analysis methods & M2.1 (architecture), M2.6--2.13, S8--S9, S13 (recalibration, locked predictions) \\
13 & Class imbalance & None used; S12.3 \\
14 & Fairness & Subgroup audit S15; limits M4 \\
15 & Model output & Probabilities, audit record and grounded report (M2.12--2.13, S15, S19); decision curves S12.7; no threshold chosen \\
16 & Development vs evaluation differences & M2.11, S13 \\
17 & Ethics & Declarations \flag{[PENDING]} \\
18a--f & Funding, conflicts, protocol, registration, data, code & Declarations \flag{[PENDING]}; plans S1; code S16 \\
19 & Patient and public involvement & Declarations (none) \\
20a--c & Flow, characteristics, comparison with development data & Table 1, S2 (Figure S1, Table S2) \\
21 & Participants and events per analysis & Tables 1, S4, S15; each result table \\
22 & Full model specification & Locked model files and model card (S15) \flag{[PENDING: release coefficients]} \\
23a--b & Performance with CIs, subgroups, heterogeneity & M3, Tables 2--3, S11--S15; heterogeneity only MU vs LUMIERE \\
24 & Model updating & S13 (recalibration, fine-tuning) \\
25--26 & Interpretation, limitations & M4 \\
27a--c & Usability, input quality, next steps & S15 (OOD flags, as-of checks), M4 \\
\bottomrule
\end{tabularx}
\end{table}

\Needspace{22\baselineskip}
\textbf{CLEAR (selected items; full table in \code{CHECKLIST\_CLEAR.md}).}

\begin{table}[H]
\centering
\footnotesize
\begin{tabularx}{\linewidth}{@{}R{1.1cm} R{5.4cm} L@{}}
\toprule
Items & Topic & Where reported \\
\midrule
1--6 & Title, abstract, keywords, background, rationale, objectives & Title, Abstract, Keywords, M1 \\
7--11 & Guidelines, ethics, sample size, study nature, eligibility & M2.2, Declarations, S12.7, M2.4, S2 \\
13--18 & Data source, overlap, split, imaging protocol, non-radiomic predictors, reference standard & M2.3, S6 (BraTS pretraining overlap risk), M2.7--2.8, S4, S3; MU protocol not distributed \\
19--24 & Segmentation, pre-processing, resampling, discretisation, image types & M2.3, S5, S14 \\
25--28 & Extraction method, classes, number, defaults & S5 \\
29--33 & Missing data, imbalance, reliability, scaling, dimension reduction & S4, S12.3, S5 (ICC filter), S5 \\
34--37 & Algorithm, tuning, confounders, model selection & M2.8, S5, S12.4 (schedule), M2.9 (phase of care) \\
38--43 & Testing, metrics, uncertainty, comparisons, explainability & M2.8, M2.11, S9, S12.6 \\
44--48 & Characteristics, flow, feature statistics, performance, comparison & Table 1, S2, S5, M3, Table 3 \\
49--52 & Discussion items & M4 \\
53--58 & Open science & Declarations, S16 \flag{[PENDING: release]} \\
\bottomrule
\end{tabularx}
\end{table}

\FloatBarrier
\section{CFB-GBM diffusion analysis}\label{sec:S18}

Source: \code{docs/results/CFB\_ADC.md};
\code{outputs/results/analyses/cfb\_adc/\{results.csv,
rule.json\}};
dataset facts \code{docs/literature/DATASETS.md}.

\textbf{Data.} TCIA CFB-GBM clinical v03, treatment and MRI availability v02, RANO v04, PyRadiomics features
v03 on the GTV; no images used. t0 is the RT-planning MRI; t1 lies a median 14 weeks after t0, about 6--8
weeks after the end of RT, inside the 12-week pseudoprogression window. t1$\rightarrow$t2 labels: 30
progressions among 107 labelled patients. No IDH, MGMT or progression dates. The RANO label is computed from
the GTV volume (RANO sizes equal the PyRadiomics GTV mesh volume, Spearman 1.00).

\textbf{Design.} One row per patient; outcome t1$\rightarrow$t2 progressive disease. Feature sets (a) clinical
(age, sex, WHO performance status, WHO edition, RT dose, fractions, delay to RT, weeks t0$\rightarrow$t1); (b)
+ log GTV volume at t1 and change t0$\rightarrow$t1; (c) + T1-Gd and FLAIR first-order and texture at t1 plus
GTV shape; (d) + ADC first-order and texture at t1. Each radiomics block: in-fold 1\% clipping, median
imputation, variance filter, scaling, top-k univariate F (k $\in$ \{5, 10\}); ridge with C $\in$ \{0.01
\ldots{} 3\}; k and C by inner stratified 3-fold CV. 20 repeats of stratified 5-fold CV, identical splits for
all sets; patient bootstrap CIs (1,000); 1,000 permutations of the ADC block for the primary contrast.

\begin{table}[H]
\centering
\small
\renewcommand{\thetable}{S25}\renewcommand{\theHtable}{S25}
\caption{\textbf{Results} (AUC = mean over 20 repeats of pooled-OOF ROC AUC).}
\label{tab:S25}
\begin{tabularx}{\linewidth}{@{}L R{3.60cm} R{3.42cm}@{}}
\toprule
Set & Conventional cohort (107 patients, 30 events) & ADC cohort (77 patients, 19 events) \\
\midrule
(a) clinical & 0.501 (0.403--0.596) & 0.565 (0.437--0.689) \\
(b) + GTV volume and change & 0.685 (0.569--0.781) & 0.713 (0.589--0.814) \\
(c) + conventional radiomics & 0.676 (0.575--0.769) & 0.691 (0.576--0.788) \\
(d) + ADC radiomics & 0.696 (0.593--0.785)\textsuperscript{a} & 0.698 (0.577--0.800) \\
\bottomrule
\end{tabularx}
\end{table}

\begin{table}[H]
\centering
\small
\begin{tabularx}{\linewidth}{@{}R{1.85cm} R{4.12cm} L@{}}
\toprule
Contrast & Conventional cohort & ADC cohort \\
\midrule
(b) vs (a) & +0.185 (+0.064, +0.293) & +0.148 (+0.031, +0.266) \\
(c) vs (b) & $-$0.009 ($-$0.067, +0.051) & $-$0.022 ($-$0.074, +0.025) \\
(d) vs (c) & +0.019 ($-$0.006, +0.045) & \textbf{+0.007 ($-$0.019, +0.035)} (primary; permutation p = 0.151; minimal detectable $\Delta$ 0.038) \\
\bottomrule
\end{tabularx}
\end{table}

\textsuperscript{a} ADC block imputed for 30 patients without ADC at t1; secondary only. \textbf{Rule not
met.} Clinical features are at chance (no molecular markers and no phase variation: every t1 scan is at the
same post-RT stage). Volume is the only informative block and is partly circular with the volumetric label.
Conventional texture adds nothing beyond volume; ADC adds nothing measurable beyond conventional radiomics.
Gains larger than about 0.04 are unlikely; smaller gains cannot be excluded. Perfusion was not available.

\FloatBarrier
\section{PRISM report generator: design and faithfulness evaluation}\label{sec:S19}

Source: \code{docs/results/REPORT\_AGENT.md};
\code{outputs/results/analyses/report\_agent/}
(\code{summary\_eval.csv}, \code{checks\_eval.csv}, \code{by\_phase.csv}, \code{stress.csv},
\code{summary\_stress.csv}, \code{reports/eval/}). Code: \code{prism/report/},
\code{prism/analyses/report\_agent.py} (\code{python -m prism report}).

\textbf{Plan.} Written after prompt and checker development on a separate development set and before any
evaluation or stress report was generated; frozen 2026-10-15 with a timestamp and the sha256 of the plan and of
every report-generator code file (\code{PLAN\_snapshot.\{md,timestamp,code\_sha256\}}). After the run only the
example-page flag labels and a duplicated ``Limitations:'' prefix were changed (figures re-run; no check, prompt
or number changed).

\textbf{Report record.} The audit record (S15) plus two blocks: a risk tier (tertile cut points of the fold
model's predicted risk on its own training scans; descriptive, not a threshold) and past-only indicative volumes
(ET and SNFH volume, change since the prior scan and the nadir, and the volumetric rule of Supplementary S12,
development AUC 0.551, labelled ``indicative, not RANO''; percent changes only from a baseline $\geq$ 1 cc). The
120-day outcome is not in the record. A flat \emph{facts} view (field key $\rightarrow$ pre-rounded value) is the only
input of the language model and the only reference of the checker. The refitted fold models reproduce the audit
stage's stored risk (maximum $|\Delta|$ \num{9.7e-17}).

\textbf{Sections.} Deterministic: scan context; inputs (value, unit, source column, as-of day); risk with 95\%
interval, tier and 80\% conformal set; contributions (intercept + clinical / timing / treatment = logit, top terms;
imaging ``not part of this model''); referral flags with reasons; volumetric change (indicative); a standing
limitations line. Language model: the summary. \emph{Kept from PRISM v1:} the change-from-prior and nadir table (now
measured volumes, past-only), the explanation section (exact additive contributions instead of SHAP on the fusion
MLP), tiered confidence language, the decision-support disclaimer. \emph{Dropped:} the fixed 0.48 alert threshold
and 0.75 ``high confidence'' cut (hand-set on the initial split; the core model's net benefit exceeds treat-all
only above about a 20\% threshold); the ``RANO 2.0 rule engine'' and RANO--AI concordance (computed from the model's
own probability, not measurements; the volumetric rule has AUC 0.55 for 120-day progression and cannot arbitrate
the model); the weighted pseudoprogression confounder score (literature weights never fitted; steroid change not in
the data; its main component is already a model input and a flag); the hard-coded molecular table and the
printed outcome label.

\textbf{Language model and arms.} Arm A (grounded): \code{Qwen/Qwen3.8-27B}, revision
\texttt{\small 1d4bf0f2\allowbreak ff6012fd\allowbreak 82039f2f\allowbreak a52739d0\allowbreak dd7c60c0} (Apache-2.0), bf16, vLLM 0.31.0 in a separate environment, tensor
parallel over 4 H100, offline (no external API; data stay on the node), non-thinking mode, greedy decoding, seed 0,
structured output \code{\{"sentences": [\{"text", "fields"\}]\}}, at most 1,024 new tokens. Frozen 10-rule system
prompt: facts only; every number copied from a field; cite field keys per sentence; state risk, interval and tier;
state every active flag and the clinician-review action, never describe an inactive flag as raised; no diagnosis
or RANO category; no management; hedged language; volumes indicative; 4--7 sentences. Arm B (original style): the
v1 \code{generate\_llm\_summary} prompt verbatim, its context blocks filled from the same record (risk in place of
the fusion probability with the original 0.48 / 0.75 phrases; the volumetric rule in the ``RANO 2.0'' block; the
original confounder score; top 3 contributions as ``SHAP drivers''; a referral-flag block added so that both arms
see the same facts); no citation constraint; same model and decoding; at most 500 new tokens. Arm C: arm A's
prompt with \code{Qwen/Qwen3.5-9B}, revision \texttt{\small c2022362\allowbreak 35762e1c\allowbreak 871ad0cc\allowbreak b60c8ee5\allowbreak ba337b9a}, 1 GPU.

\textbf{Checker} (deterministic, patterns frozen with the plan). Per claim: every cited key must exist; every number
must equal a number in a cited field up to the written rounding (magnitudes of signed values and $\times$100 for
percentages allowed); a number without a citation is unsupported. Arm B has no citations; its numbers are checked
against its own prompt (``vs own input'') and against the record. Per report: each active flag mentioned (keyword
patterns per flag); inactive flags not described as raised; tier words consistent with the record; no forbidden
content (diagnosis, RANO category as outcome, management, unhedged certainty; full patterns in the source
document). An unparsed report fails every check. The checker cannot see a correct number attached to the wrong
concept; meaning was checked by an AI assistant's reading (below).

\textbf{Sets.} Development: the 59 non-grade-4 scans (two rounds; final pass A 58/59, B 51/59, C 54/59).
Evaluation: all 249 grade 4 scans (126 patients), records from the out-of-fold fits of CV repeat 0, one report per
scan per arm. Primary (pre-specified): arm A, share of reports with zero unsupported claims and every active flag
mentioned, Wilson 95\% CI, target $\geq$ 0.95. Secondary: paired differences on the same scans, 2,000-draw
patient-cluster bootstrap (seed 0). Stress: 30 perturbed evaluation records (10 risk moved to another tier's
mid-point, 10 flag switched on, 10 flag switched off); a report passes if it passes the checks against the
perturbed record and follows the change.

\begin{table}[H]
\centering
\footnotesize
\renewcommand{\thetable}{S26}\renewcommand{\theHtable}{S26}
\caption{\textbf{Report faithfulness on 249 grade 4 scans} (\code{summary\_eval.csv}). CI: Wilson 95\%.}
\label{tab:S26}
\begin{tabularx}{\linewidth}{@{}L R{3.4cm} R{3.6cm} R{3.4cm}@{}}
\toprule
 & A grounded, 27B & B original prompt, 27B & C grounded, 9B \\
\midrule
Pass (zero unsupported + all flags) & 236/249 (0.948) [0.913, 0.969] & 171/249 (0.687) [0.627, 0.741] (vs own input) & 243/249 (0.976) [0.948, 0.989] \\
Pass vs the record & 236 & 0 & 243 \\
Reports with an unsupported claim & 13 & 38 (own input); 249 (record) & 6 \\
Unsupported claims per report & 0.05 & 0.15 (own input); 2.33 (record) & 0.02 \\
Active flags omitted & 0/269 & 48/269 (17.8\%), 47 reports & 0/269 \\
Spurious flag & 0 & 1 & 0 \\
Forbidden content, any & 0\% & 100\% & 2.8\% (7) \\
\quad diagnosis / RANO category / management / certainty & 0 / 0 / 0 / 0\% & 51\% / 41\% / 99\% / 40\% & 0.8\% / 0 / 2.0\% / 0.8\% \\
Tier stated / inconsistent & 100\% / 0\% & 57\% / 44\% (other scheme) & 100\% / 0.4\% \\
Batched time per report & 0.06 s (4 H100) & 0.06 s & 0.08 s (1 H100) \\
Single-request latency, median of 10 & 3.6 s & -- & 3.1 s \\
\bottomrule
\end{tabularx}
\end{table}

\textbf{Arm A.} All 13 failures: the risk sentence says ``within 120 days'' without citing \code{risk.horizon\_days};
the number is correct and in the record (0 numbers absent from the record). All 13 are early post-RT scans (pass by
phase: pre-RT 76/76, during RT 6/6, early post-RT 85/98, later post-RT 69/69); the development set showed the same
mode once. 249/249 parsed; 0 invalid citations; 0 truncated; 5.9 sentences and 480 output tokens per report.
\emph{EXPLORATORY:} accepting a number that appears verbatim in a field cited elsewhere in the report would give
249/249; this relaxation is post hoc.

\textbf{A vs B} (paired, 249 scans): pass +0.261 [+0.197, +0.328]; pass vs record +0.948 [+0.919, +0.973];
unsupported claims per report $-$0.10 [$-$0.15, $-$0.05]; flag omissions per report (keyword-based for both)
$-$0.19 [$-$0.25, $-$0.13]; forbidden content $-$1.00 [$-$1.00, $-$1.00]. Arm B states a RANO verdict in 241 reports,
mentions pseudoprogression in 249, recommends management in 242 (MDT / tumour board 166, short-interval MRI 172),
mentions steroids (not in the record) in 84 and calls the 120-day horizon ``12-month'' in 41. \textbf{A vs C:} pass
$-$0.028 [$-$0.062, +0.008]; forbidden $-$0.028 [$-$0.053, $-$0.008] (C worse). Of C's 7 forbidden hits, 1 is a real
violation (``suggesting changes that may reflect treatment effects'') and 6 are pattern false positives (for example
``this does not confirm progression''); one C report also gives a wrong reason for a referral, a meaning error the
checker cannot see. C's 6 unsupported claims are 4 invalid keys and 2 citation omissions.

\begin{table}[H]
\centering
\small
\renewcommand{\thetable}{S27}\renewcommand{\theHtable}{S27}
\caption{\textbf{Stress test: 30 perturbed records} (\code{stress.csv}; pass = checks against the perturbed record and
the report follows the change; checks alone: A 28, B 17, C 29).}
\label{tab:S27}
\begin{tabularx}{\linewidth}{@{}L R{2.6cm} R{2.4cm} R{2.4cm} R{1.4cm}@{}}
\toprule
Arm & Risk moved to another tier & Flag switched on & Flag switched off & Total \\
\midrule
A & 8/10 & 10/10 & 9/10 & 27/30 \\
B & 2/10 & 6/10 & 5/10 & 17/30 \\
C & 9/10 & 9/10 & 10/10 & 28/30 \\
\bottomrule
\end{tabularx}
\end{table}

The three arm-A failures, read: the horizon-citation omission (the new risk and tier were stated correctly); a new
interval bound that coincided with the old risk (the new risk, interval and tier were stated correctly; C fails
the same record); and an invented key (\code{risk.tier\_cut\_upper\_percent}) for a value that is in the record.
\emph{EXPLORATORY:} in all 30 perturbed records the grounded text followed the changed field.

\textbf{Reading by an AI assistant (not a clinician).} The three pre-specified examples (first scan per phase) and 12
random arm-A reports (seed 1) were read against their facts by an AI assistant (Claude); no human or clinician has
read them yet: no factual error; every active flag stated with its reason. Wording issues the
checker cannot see: ``95\% confidence interval'' for a bootstrap interval (twice); a double negative for a volume
below the nadir; a pre-diagnosis scan whose ``all inputs known by the scan day'' is ``no'' (the molecular look-ahead,
S15) is not mentioned in the summary, only in the deterministic section; a risk just below a tier cut point that
rounds to the same value as the cut. Arm B on the same examples states ``stable disease'' or ``Progressive Disease
(PD)'', recommends MDT review and short-interval MRI, and mentions steroid dosing that is not in the record.

\textbf{Open issues.} The primary target was missed by one report; the obvious fix (require citing the horizon
field, or treat constants as citable) would be tuned on the evaluation set and needs fresh records (for example
another CV repeat). The checker is lexical and has false positives on negations. Tier cut points should be shown
with more decimals. Re-run identity of greedy decoding under vLLM batching was not tested. Records come from one CV
repeat of one cohort, and the core model has no external validation with treatment fields. No clinician evaluated
usefulness, readability or trust.

\FloatBarrier
\section*{Supplementary figures}
\addcontentsline{toc}{section}{Supplementary figures}

\begin{itemize}
\item \textbf{Figure S1.} Cohort flow (\code{docs/manuscript/cohort\_flow.png}).
\item \textbf{Figure S2.} Split-variance histograms per model, original analysis figure
(\code{outputs/results/analyses/split\_variance/split\_variance.png};
pre-clip radiomics arm --- superseded by
main Figure 2).
\item \textbf{Figure S3.} Phase-ladder overview
(\code{outputs/results/analyses/phase\_ladder/ladder.png}).
\item \textbf{Figure S4.} Sensitivity forest plot
(\code{outputs/results/analyses/sensitivity/forest.png}).
\item \textbf{Figure S5.} Calibration curves and decision curves
(\code{outputs/results/analyses/utility/calibration\_dca.png}).
\item \textbf{Figure S6.} Imaging diagnostics: label--date alignment and horizon curve
(\code{outputs/results/analyses/diagnostics\_imaging/\{alignment,horizon\_curve\}.png});
visual QA panels
(\code{qa/*.png}).
\item \textbf{Figure S7.} SHAP summaries for the baseline and full models
(\code{outputs/results/analyses/shap/grade4/*\_\{bar,beeswarm\}.png}).
\item \flag{[PENDING: export to the journal format.]}
\end{itemize}

\FloatBarrier
\bibliography{refs}

\end{document}
